%% 8-9 わき出しと吸いこみ：発散 div（第14回：組 V8-5。古い見本 ch8-9.tex の図・池のデータを生かして書き直した。ページ制限なし。偶数ページ）
%% ほかの巻・物理の本との約束：
%%   物理 p1-13:388・392「その点の密度が増える速さ＋その点からのわき出し＝0（数学 8-9 の発散）」、div v = ∂vx/∂x+∂vy/∂y+∂vz/∂z = 0。
%%   8-8 の流束（∯・\vec n dS・場合G：閉じた箱を面ごとに）→ この章の「せいかくに言うと」の定義と例題6。
%%   池（わき出し (1.5, 2.5)、吸いこみ (4.5, 1.5)、Q=4 m^2/s）は 8-5・8-11 と同じ数。主役の場（平面の6つ）は 8-5・8-10・8-11 と同じ。
%%   立方体 (x^2, xy, -z) の div 3x-1 は 8-8（面ごとに 0.5）・8-11（∭ と一致）と同じ場。
%%   点のわき出し（池の口・点電荷）は「8-11 でくりぬく」へ送る。
\makeatletter
% ---- 「名前と操作」の型（tools/kata8-int.md。kviiig → kviiij）----
\colorlet{kviiijNc}{bkOrange!85!black}% 名前の枠の線（点線）
\colorlet{kviiijSc}{bkTheme}% ふつうの定積分の枠の線（実線）
\colorlet{kviiijSb}{bkTheme!8}% ふつうの定積分の地
\tikzset{
  kviiij dan/.style={font=\sffamily\bfseries\footnotesize, rounded corners=1.2mm, inner xsep=1.6mm, inner ysep=0.5mm, line width=0.9pt},
  kviiij dan1/.style={kviiij dan, draw=kviiijNc, dash pattern=on 2.2pt off 1.4pt, fill=white, text=bkInk},
  kviiij dan2/.style={kviiij dan, draw=kviiijSc, fill=kviiijSb, text=bkInk},
  kviiij dan3/.style={kviiij dan, draw=bkTheme, fill=bkTheme, text=bkOn},
  kviiij name/.style={draw=kviiijNc, line width=0.9pt, dash pattern=on 2.2pt off 1.4pt, rounded corners=1.2mm, inner xsep=1.2mm, inner ysep=1.2mm, fill=white},
  kviiij seki/.style={draw=kviiijSc, line width=0.9pt, rounded corners=1.2mm, inner xsep=1.2mm, inner ysep=1.2mm, fill=kviiijSb},
  kviiij tag/.style={font=\sffamily\bfseries\footnotesize, text=bkInk, inner sep=0.4mm},
  kviiij q/.style={diamond, aspect=2.4, draw=bkYellow!75!black, line width=0.9pt, fill=bkYellowBg, font=\sffamily\footnotesize, align=center, inner sep=0.3mm, text=bkInk},
  kviiij do/.style={rectangle, draw=\bkTT, line width=0.9pt, fill=white, rounded corners=0.6mm, font=\sffamily\footnotesize, align=center, inner xsep=1.5mm, inner ysep=1.2mm, text=bkInk},
  kviiij go/.style={-{Stealth[length=2.2mm]}, line width=0.9pt, bkInk!75, shorten <=1.5mm, shorten >=1.5mm},
  kviiij yn/.style={font=\sffamily\bfseries\footnotesize, text=bkInk, fill=white, inner sep=0.5mm},
}
% 段の札（例題の中で）  \kviiijdan{1|2|3}
\newcommand\kviiijdan[1]{\tikz[baseline=(d.base)]\node[kviiij dan#1] (d) {\vphantom{あ}\ifcase#1\or 名前\or 直す\or 楽をする\fi};}
% 段の見出し（本文で）  \kviiijdanhead{1|2|3}{見出し}
\newcommand\kviiijdanhead[2]{\par\addvspace{8pt}\Needspace{7\baselineskip}\noindent
  \tikz[baseline=(d.base)]\node[kviiij dan#1, font=\sffamily\bfseries\small, inner ysep=0.8mm] (d) {\vphantom{あ}\ifcase#1\or ①\,名前の段\or ②\,直す段\or ③\,楽をする段\fi};%
  \hspace{0.7\zw}{\sffamily\bfseries\color{\bkTT}#2}\par\nobreak\vspace{3pt}\noindent\ignorespaces}
% 場合の札  \kviiijcase{C}（′ は \kviiijcase{A$'$}）
\newcommand\kviiijcase[1]{\tikz[baseline=(c.base)]\node[fill=bkTheme, text=bkOn, font=\sffamily\bfseries\small, rounded corners=0.8mm, inner xsep=1mm, inner ysep=0.6mm, minimum width=5.4mm] (c) {\vphantom{Ag}#1};}
% 例題の1行め  \kviiijbaai{ことば}{8-6}{C}（場合が2つなら {C,G}）
\newcommand\kviiijbaai[3]{{\sffamily\bfseries 場合：}#1（表#2の\kviiij@cases#3,\@nil）}
\def\kviiij@cases#1,#2\@nil{\kviiijcase{#1}\ifx\relax#2\relax\else と\kviiij@cases#2\@nil\fi}
% 式の中の「名前」と「ふつうの定積分」
\newcommand\kviiijN[1]{\tikz[baseline=(n.base)]\node[kviiij name] (n) {$\displaystyle #1$};}
\newcommand\kviiijS[1]{\tikz[baseline=(n.base)]\node[kviiij seki] (n) {$\displaystyle #1$};}
\newcommand\kviiijNt[1]{\tikz[baseline=(n.base)]{\node[kviiij name] (n) {$\displaystyle #1$};\node[kviiij tag, anchor=north] at ([yshift=-0.8mm]n.south) {名前};}}
\newcommand\kviiijSt[1]{\tikz[baseline=(n.base)]{\node[kviiij seki] (n) {$\displaystyle #1$};\node[kviiij tag, anchor=north] at ([yshift=-0.8mm]n.south) {ふつうの定積分};}}
% 場合分けの表の題の帯  \kviiijcasetitle{8-6}{題}
\newcommand\kviiijcasetitle[2]{\par\addvspace{6pt}\Needspace{8\baselineskip}\noindent\begin{tikzpicture}[baseline=(t.base)]
  \node[fill=bkTheme, text=bkOn, font=\sffamily\bfseries\small, rounded corners=1mm, inner xsep=2mm, inner ysep=1mm] (t) {場合分けの表\,#1};
  \node[anchor=west, font=\sffamily\bfseries, text=\bkTT, inner sep=0pt] at ([xshift=2mm]t.east) {#2};
  \end{tikzpicture}\par\nobreak\vspace{3pt}}
% ---- この章の部品 ----
\tikzset{
  kviiij arrow/.style={-{Stealth[length=3.2pt, width=3pt]}, line width=0.9pt, bkTheme},
  kviiij pond/.style={-{Stealth[length=3.4pt, width=3.2pt]}, line width=0.9pt, bkTheme},
  kviiij box/.style={bkOrange!90!black, line width=1.2pt, dash pattern=on 3pt off 2pt},
  kviiij big/.style={-{Stealth[length=6pt, width=5pt]}, line width=1.4pt, bkTheme},
}
% ベクトル場の小さな図  \kviiijpanel{vx の式}{vy の式}{図の中に足すもの}（\x, \y を使う。-2〜2 の格子、1目もり 5.6 mm）
\newcommand\kviiijpanel[3]{%
  \begin{tikzpicture}[x=5.6mm, y=5.6mm]
    \fill[white] (-2.6,-2.6) rectangle (2.6,2.6);
    \draw[bkLine, line width=0.6pt] (-2.6,-2.6) rectangle (2.6,2.6);
    \foreach \x in {-2,...,2} {\foreach \y in {-2,...,2} {
      \pgfmathsetmacro\kviiiju{#1}\pgfmathsetmacro\kviiijw{#2}
      \pgfmathtruncatemacro\kviiijok{(abs(\kviiiju)+abs(\kviiijw))>0.02}
      \ifnum\kviiijok=1 \draw[kviiij arrow] ({\x-\kviiiju/2},{\y-\kviiijw/2}) -- ++(\kviiiju,\kviiijw);\fi}}
    #3
  \end{tikzpicture}}
% 池の矢印（6 m×4 m。格子点の速さの向き。長さは速さに合わせて縮めてある：速いほど長い）
\newcommand\kviiijpondarrows{%
  \draw[kviiij pond] (-0.011,0.106) -- ++(0.022,-0.213);
  \draw[kviiij pond] (0.936,0.110) -- ++(0.129,-0.220);
  \draw[kviiij pond] (1.873,0.071) -- ++(0.254,-0.143);
  \draw[kviiij pond] (2.841,-0.012) -- ++(0.319,0.025);
  \draw[kviiij pond] (3.881,-0.119) -- ++(0.238,0.238);
  \draw[kviiij pond] (5.004,-0.152) -- ++(-0.007,0.305);
  \draw[kviiij pond] (6.068,-0.099) -- ++(-0.136,0.198);
  \draw[kviiij pond] (0.044,1.120) -- ++(-0.089,-0.241);
  \draw[kviiij pond] (0.976,1.166) -- ++(0.048,-0.333);
  \draw[kviiij pond] (1.853,1.132) -- ++(0.295,-0.264);
  \draw[kviiij pond] (2.785,1.031) -- ++(0.429,-0.061);
  \draw[kviiij pond] (3.767,0.851) -- ++(0.467,0.297);
  \draw[kviiij pond] (5.156,0.815) -- ++(-0.313,0.370);
  \draw[kviiij pond] (6.138,0.954) -- ++(-0.275,0.092);
  \draw[kviiij pond] (0.128,2.075) -- ++(-0.256,-0.151);
  \draw[kviiij pond] (1.143,2.207) -- ++(-0.286,-0.413);
  \draw[kviiij pond] (1.766,2.182) -- ++(0.467,-0.363);
  \draw[kviiij pond] (2.762,2.079) -- ++(0.477,-0.159);
  \draw[kviiij pond] (3.766,2.182) -- ++(0.467,-0.363);
  \draw[kviiij pond] (5.143,2.207) -- ++(-0.286,-0.413);
  \draw[kviiij pond] (6.128,2.075) -- ++(-0.256,-0.151);
  \draw[kviiij pond] (0.138,2.954) -- ++(-0.275,0.092);
  \draw[kviiij pond] (1.156,2.815) -- ++(-0.313,0.370);
  \draw[kviiij pond] (1.767,2.851) -- ++(0.467,0.297);
  \draw[kviiij pond] (2.785,3.031) -- ++(0.429,-0.061);
  \draw[kviiij pond] (3.853,3.132) -- ++(0.295,-0.264);
  \draw[kviiij pond] (4.976,3.166) -- ++(0.048,-0.333);
  \draw[kviiij pond] (6.044,3.120) -- ++(-0.089,-0.241);
  \draw[kviiij pond] (0.068,3.901) -- ++(-0.136,0.198);
  \draw[kviiij pond] (1.004,3.848) -- ++(-0.007,0.305);
  \draw[kviiij pond] (1.881,3.881) -- ++(0.238,0.238);
  \draw[kviiij pond] (2.841,3.988) -- ++(0.319,0.025);
  \draw[kviiij pond] (3.873,4.071) -- ++(0.254,-0.143);
  \draw[kviiij pond] (4.936,4.110) -- ++(0.129,-0.220);
  \draw[kviiij pond] (5.989,4.106) -- ++(0.022,-0.213);}
% 池の格子
\newcommand\kviiijpondgrid{%
  \foreach \i in {1,...,5} {\draw[bkLine, line width=0.5pt] (\i,0) -- (\i,4);}
  \foreach \j in {1,...,3} {\draw[bkLine, line width=0.5pt] (0,\j) -- (6,\j);}
  \draw[bkGray, line width=0.8pt] (0,0) rectangle (6,4);}
% 発散の地図のマス  \kviiijcell{列}{行（下から）}{値の文字}{0:白 1:わき出し 2:吸いこみ}
\newcommand\kviiijcell[4]{%
  \ifcase#4\relax\or\fill[bkOrange!35] (#1,#2) rectangle ++(1,1);\or\fill[bkBlue!25] (#1,#2) rectangle ++(1,1);\fi
  \node[font=\sffamily\small\bfseries, text=bkInk] at (#1+0.5,#2+0.5) {#3};}
\newcommand\kviiijpondmap{%
  \foreach \i in {0,...,5} {\foreach \j in {0,...,3} {\kviiijcell{\i}{\j}{0.0}{0}}}
  \fill[white] (1.03,2.03) rectangle (1.97,2.97);\kviiijcell{1}{2}{4.0}{1}
  \fill[white] (4.03,1.03) rectangle (4.97,1.97);\kviiijcell{4}{1}{\raisebox{0.3em}{\rule{0.45em}{0.09em}}\hspace{0.06em}4.0}{2}}
% 小見出し（太字の1行）
\newcommand\kviiijhead[1]{\par\addvspace{6pt}\Needspace{6\baselineskip}\noindent{\sffamily\bfseries\color{\bkTT}#1}\par\nobreak\vspace{2pt}}
% 囲みの中の小見出し（それだけで1行）
\newcommand\kviiijbox[1]{\par\addvspace{3pt}\noindent\textbf{#1}\par\vspace{3pt}\noindent\ignorespaces}
% 式を、まん中に1行で置く
\newcommand\kviiijline[1]{\par\addvspace{3pt}{\centering #1\par}\addvspace{3pt}\noindent\ignorespaces}
% 大事な式のわく  \kviiijframe[はば]{中身}
\newcommand\kviiijframe[2][100mm]{\par\addvspace{4pt}{\centering\tikz\node[draw=\bkTT, line width=0.9pt, fill=bkTheme!10, rounded corners=1.5mm, inner xsep=6mm, inner ysep=2.5mm, align=center]{#2};\par}\addvspace{4pt}\noindent\ignorespaces}
% 例題の途中式の2行め以降
\newcommand\kviiija{\par\vspace{2pt}\noindent\rightskip=0pt plus 8em\relax\hangindent=2.6\zw\hangafter=0\relax}
\newcommand\kviiijb{\par\vspace{2pt}\noindent\rightskip=0pt plus 8em\relax\hangindent=4.3\zw\hangafter=0\relax}
\newcommand\kviiijs{\rightskip=0pt plus 8em\relax}
\newcommand\kviiijpd[2]{\frac{\partial #1}{\partial #2}}
\newcommand\kviiijdiv{\mathrm{div}\,}
\newcommand\kviiijps{\ensuremath{\mathrm{/s}}}% 「毎秒」（\unit{/s} は / が消えるので使わない）
\makeatother

\bkchapter{8}{9}{わき出しと吸いこみ：\\発散 div}{\bklevel{5}}{わき出しと吸いこみ：発散 div}
\bkmichishirube{4-12,4-13,8-2,8-4,8-5,8-8}{8-10,8-11,8-12,8-13,8-14}

\begin{goal}
\begin{itemize}[label={\color{\bkTT}●}, leftmargin=1.6\zw, itemsep=0pt, topsep=1pt]
\item \bklevel{3}\enspace 矢印の図にマスを置いて、わき出しと吸いこみを見分けられる。
\item \bklevel{4}\enspace マスの辺を通る量から、表計算で発散を求められる。
\item \bklevel{5}\enspace 発散を偏微分の和で計算し、流束や連続の式と結びつけられる。
\end{itemize}
\end{goal}

\begin{hatena}
広い池の水面の一部（\qty{6}{m}×\qty{4}{m}）で、水の流れを調べました。矢印は、その点の水の流れの向きと速さです（速いほど長くかいてあります）。池の底には、水がわき出してくる口と、水が吸いこまれる口が、1つずつかくれています。どこにあるでしょう。
\begin{bkfigure}
\begin{tikzpicture}[x=9mm, y=9mm]
  \fill[bkBlueBg, rounded corners=2mm] (-0.45,-0.45) rectangle (6.45,4.45);
  \kviiijpondarrows
  \foreach \i in {0,...,6} {\draw[bkGray, line width=0.6pt] (\i,-0.45) -- (\i,-0.6); \node[bk tick, below=0.5mm] at (\i,-0.6) {\i};}
  \foreach \j in {0,...,4} {\draw[bkGray, line width=0.6pt] (-0.45,\j) -- (-0.6,\j); \node[bk tick, left=0.5mm] at (-0.6,\j) {\j};}
  \node[bk label, anchor=west] at (6.6,-0.85) {$x$（m）};
  \node[bk label, anchor=south west] at (-0.4,4.55) {$y$（m）};
\end{tikzpicture}
\end{bkfigure}
\end{hatena}

\noindent 8-5 では、場所ごとに矢印が1本ずつある「ベクトル場」を、図にかいて見ました。8-8 では、矢印の場が面を通りぬける量（流束）を、面積分で求めました。この章では、その2つを組み合わせて、「\textbf{水がどこで生まれて、どこで消えているか}」を見つける道具を作ります。

矢印をながめているだけでは、わき出しの口はなかなか見つかりません。そこで、流れの中に小さな「箱」を置き、箱から出ていく水と入ってくる水を数えます。数えるだけなら、たし算とひき算でできます。最後に、箱を小さくしていくと、8-2 の偏微分が出てきます。

★★★★★の章ですが、「目でわかる」「手でわかる」の段は、図と表計算だけで読めるように書きました。「式でわかる」の段は、1回で分からなくても大丈夫です。図と数にもどりながら、ゆっくり進みましょう。

%% ---------------- やってみよう・思い出そう ----------------
\begin{yattemiyou}[水を出す所、吸いこむ所を見よう]
\textbf{用意するもの}\quad 浅いバット（なければ大きな平たい皿）、水、こしょう（またはちぎった紙くず）、スポイト2本
\begin{enumerate}[label={\bkcircled[bkGreen]{\arabic*}}, leftmargin=2.4\zw, itemsep=1pt, topsep=3pt]
\item バットに深さ\qty{1}{cm}ほど水を入れ、こしょうを水面にうすくふる。
\item まん中で、スポイトから水をゆっくり出す。こしょうの動きを、ノートに矢印でスケッチしよう。
\item 水が静かになったら、こんどはスポイトで水をゆっくり吸いこむ。同じようにスケッチしよう。
\item 2本のスポイトで、出すのと吸うのを同時にやってみよう。
\end{enumerate}
\begin{bkfigure}
\begin{tikzpicture}[x=1mm, y=1mm]
  % 横から見たところ：バット・水・スポイト
  \fill[bkBlue!18] (2,2) rectangle (52,8);
  \draw[bkGray, line width=1pt] (0,14) -- (2,2) -- (52,2) -- (54,14);
  \foreach \x in {6,11,...,48} {\fill[bkInk!70] (\x,8.2) circle (0.35);}
  \draw[bkInk!80, line width=0.8pt, fill=white] (25.5,9.5) rectangle (28.5,24);
  \draw[bkInk!80, line width=0.8pt, fill=bkRed!50] (27,27.5) ellipse (2.6 and 3.6);
  \draw[bk arrow, bkBlue] (27,9) -- (27,5.5);
  \node[anchor=north, font=\sffamily\bfseries\small, text=bkInk] at (27,-1.5) {横から見たところ};
  % 上から見たところ：こしょうをうすくふった水面
  \begin{scope}[shift={(72,0)}]
  \draw[bkGray, line width=1pt, fill=bkBlue!12, rounded corners=1.5mm] (0,0) rectangle (56,30);
  \foreach \x/\y in {5/5,12/22,18/9,25/26,31/4,38/18,44/8,50/24,8/14,21/16,35/12,47/15,15/27,29/20,41/27,52/5}
    {\fill[bkInk!70] (\x,\y) circle (0.45);}
  \draw[bkRed!70, line width=0.8pt] (28,15) circle (1.6);
  \node[anchor=north, font=\sffamily\bfseries\small, text=bkInk] at (28,-1.5) {上から見たところ};
  \end{scope}
\end{tikzpicture}
\end{bkfigure}
\end{yattemiyou}

\noindent 水を出した所のまわりでは、こしょうは四方へ広がり、矢印はみな\textbf{外向き}になります。吸いこんだ所のまわりでは、矢印はみな\textbf{内向き}です。2本を同時に使うと、出した所から吸った所へ、ぐるっと回りこむ流れができます。\bkhatenaref の池とよく似た流れです。

水がわいて出てくる所を\textbf{わき出し}（わきだし）、水が消えていく所を\textbf{吸いこみ}（すいこみ）といいます。\bkidx{わきだし}{わき出し}\bkidx{すいこみ}{吸いこみ}スポイトのように口が見えていればすぐわかります。でも、池の底の口は見えません。見えるのは、水面の矢印だけです。矢印から口を見つける方法を考えましょう。

\begin{omoidasou}
\bkrecall{8-8}{面を通りぬける量（流束）は、流れの「面に直角な成分」×面積。向きは、面の法線ベクトル $\vec n$ で決める。閉じた面では $\vec n$ を外向きにとり、外へ出る量を正にする。}
\bkrecall{8-2}{偏微分 $\dfrac{\partial f}{\partial x}$ は、ほかの文字（$y$ や $z$）を止めて、$x$ だけを動かしたときの変化の割合。止めた文字は数と同じに扱う。}
\bkrecall{8-5}{主役の場：広がる $(x,\ y)$、集まる $(-x,\ -y)$、そろった流れ $(1,\ 0)$、うず $(-y,\ x)$、ずれ $(y,\ 0)$、側溝 $(x,\ 0)$。$(x,\ y)$ は場所、出てくる組 $(v_x,\ v_y)$ は矢印。}
\bkrecall{8-4}{ナブラ $\nabla=\Bigl(\dfrac{\partial}{\partial x},\ \dfrac{\partial}{\partial y},\ \dfrac{\partial}{\partial z}\Bigr)$。$\nabla f$ は $\mathrm{grad}\,f$ と同じ。}
\bkrecall{4-12}{内積 $\vec a\cdot\vec b=a_xb_x+a_yb_y+a_zb_z$。成分どうしをかけて、たす。答えは1つの数。}
\end{omoidasou}

%% ---------------- 数学の道具 ----------------
\bksection{doguu}{数学の道具}

\bkstage{me}{目でわかる}{小さな箱を置いてみる}{\bklevel{3}}

\bksubsection{1}{箱から出る量と入る量をくらべる}

\label{bk:8-9:hako}流れの中に、目に見えない小さな箱を置いてみます。平面の流れなら、箱は方眼の1マスです。マスの4つの辺を通って、1秒間に
\kviiijframe[96mm]{\sffamily\bfseries （外へ出ていく水の量）－（中へ入ってくる水の量）}
を数えます。この差を「\mbox{出－入}」とよぶことにします。3つの場合を図にしました。

\begin{bkfigure}
\begin{tikzpicture}[x=1mm, y=1mm]
  % 左：わき出しのマス
  \begin{scope}[shift={(8,0)}]
    \draw[kviiij box] (0,0) rectangle (24,24);
    \fill[bkOrange!90!black] (12,12) circle (1.3);
    \foreach \a in {0,90,180,270} {\draw[bk arrow, bkTheme, line width=1.4pt] (12,12) ++(\a:6) -- ++(\a:14);}
    \node[font=\sffamily\small, align=center, anchor=north, text=bkInk] at (12,-10.5) {出る 4　入る 0\\[1pt]{\bfseries\ltjsetparameter{autoxspacing=false}出－入＝＋4}};
  \end{scope}
  % 中：一様な流れ
  \begin{scope}[shift={(58,0)}]
    \draw[kviiij box] (0,0) rectangle (24,24);
    \foreach \y in {5,12,19} {\draw[bk arrow, bkTheme, line width=1.4pt] (-10,\y) -- (-1,\y);}
    \foreach \y in {5,12,19} {\draw[bk arrow, bkTheme, line width=1.4pt] (25,\y) -- (34,\y);}
    \node[font=\sffamily\small, align=center, anchor=north, text=bkInk] at (12,-10.5) {出る 3　入る 3\\[1pt]{\bfseries 出－入＝0}};
  \end{scope}
  % 右：吸いこみのマス
  \begin{scope}[shift={(104,0)}]
    \draw[kviiij box] (0,0) rectangle (24,24);
    \fill[bkBlue!80!black] (12,12) circle (1.3);
    \foreach \a in {0,90,180,270} {\draw[bk arrow, bkTheme, line width=1.4pt] (12,12) ++(\a:19) -- ++(\a:-12);}
    \node[font=\sffamily\small, align=center, anchor=north, text=bkInk] at (12,-10.5) {出る 0　入る 4\\[1pt]{\bfseries\ltjsetparameter{autoxspacing=false}出－入＝－4}};
  \end{scope}
\end{tikzpicture}
\bkfigcaption{点線のマスの辺を通る矢印を数える（矢印1本を「1」と数えた）。左はマスの中に水の出る口、右は吸いこむ口がある。中は口がなく、入った分だけ出ていく。}
\label{bk:8-9:fig-hako}
\end{bkfigure}

\begin{itemize}[label={\color{\bkTT}●}, leftmargin=1.6\zw, itemsep=2pt, topsep=2pt]
\item 「\mbox{出－入}」が\textbf{正}なら、マスの中で水が生まれています。マスの中は\textbf{わき出し}です。
\item 「\mbox{出－入}」が\textbf{負}なら、マスの中で水が消えています。マスの中は\textbf{吸いこみ}です。
\item 「\mbox{出－入}」が\textbf{0}なら、入った分だけ出ていきます。わき出しでも吸いこみでもありません。
\end{itemize}

なぜ、これでわかるのでしょう。水は、勝手にふえたり消えたりしません。マスの中に口がなければ、マスに入った水は、どこかの辺から必ず出ていきます。だから「\mbox{出－入}」は0です。0でないのは、マスの中に、水を足したり取ったりする口があるときだけです。

\kviiijhead{大きさをそろえる：1平方メートルあたりにする}
出入りの量は、箱の大きさでも変わります。同じ流れの中でも、大きな箱ほど、たくさんの量が出入りします。たとえば、\textbf{雨水がどこにでも少しずつ流れこむ}みぞでは、長さ\qty{1}{m}の区間の出入りの差より、\qty{2}{m}の区間の差のほうが2倍大きくなります。これでは、場所どうしをくらべられません。

そこで、「\mbox{出－入}」を箱の面積でわって、\textbf{1平方メートルあたり}の量にそろえます。人口を面積でわって「1平方キロメートルあたりの人数（人口密度）」にするのと同じ考え方です。こうすれば、箱の大きさがちがっても、同じ物差しでくらべられます。

\begin{taisetsu}
小さな箱から出ていく量（\mbox{出－入}）を、箱の大きさ（平面なら面積、空間なら体積）でわったものを、その場所の\textbf{発散}（はっさん）といいます。\bkidx{はっさん}{発散}
\kviiijline{発散 $=\dfrac{\text{出}-\text{入}}{\text{箱の大きさ}}$\qquad 記号は $\mathrm{div}\,\vec v$（ダイバージェンス）}
発散が正の所が\textbf{わき出し}、負の所が\textbf{吸いこみ}。
\end{taisetsu}

$\mathrm{div}$ は、英語の divergence（ダイバージェンス：広がって出ていくこと）の頭の3文字です。$\vec v$ は流れの速さの矢印（速度）です。「$\vec v$ という矢印の場の、発散」と読みます。

\begin{column}{「発散」という名前}
divergence ということばを、いまの意味で使ったのは、イギリスの数学者クリフォードで、1878年の本『動力学の要素（Elements of Dynamic）』だと伝えられています。それより前に、電磁気学の式を整えたマクスウェルは、符号を反対にした量を convergence（集まること：収束）とよんでいました。吸いこみの強さをはかる、という向きの名前です。日本語の「発散」は、divergence を訳したことばです。（数学には、数列が限りなく大きくなる「発散」もあります。名前は同じですが、別の意味です。）
\end{column}

\Needspace{12\baselineskip}
\bksubsection{2}{いろいろな流れでくらべる}

\label{bk:8-9:kuraberu}8-5 の主役の場6つに、同じ大きさのマス（点線）を置いて、出入りをくらべてみます。矢印は、見やすいように縮めてかいてあります。

\begin{bkfigure}
\begin{tabular}{@{}>{\centering\arraybackslash}p{40mm}@{\hspace{3mm}}>{\centering\arraybackslash}p{40mm}@{\hspace{3mm}}>{\centering\arraybackslash}p{40mm}@{}}
\kviiijpanel{0.32*\x}{0.32*\y}{\draw[kviiij box] (-0.5,-0.5) rectangle (0.5,0.5);} &
\kviiijpanel{-0.32*\x}{-0.32*\y}{\draw[kviiij box] (-0.5,-0.5) rectangle (0.5,0.5);} &
\kviiijpanel{0.6}{0}{\draw[kviiij box] (-0.5,-0.5) rectangle (0.5,0.5);} \\[1mm]
{\sffamily\bfseries\small（ア）広がる $(x,\ y)$} & {\sffamily\bfseries\small（イ）集まる $(-x,\ -y)$} & {\sffamily\bfseries\small（ウ）そろった流れ $(1,\ 0)$} \\
{\sffamily\small 出＞入：発散は正} & {\sffamily\small 出＜入：発散は負} & {\sffamily\small 出＝入：発散は0} \\
\end{tabular}
\bkfigcaption{主役の場（ア）〜（ウ）。}
\label{bk:8-9:fig-six}
\end{bkfigure}

1つずつ、マスの辺を見ていきます。
\begin{enumerate}[label={\bkcircled[\bkTT]{\small\arabic*}}, leftmargin=2.4\zw, itemsep=3pt, topsep=2pt]
\item \textbf{（ア）広がる}\par\nobreak\noindent マスのどの辺でも、矢印は外へ向いています。出る一方なので「\mbox{出－入}」は正です。マスをどこに置いても同じで、どの場所も少しずつわき出しています。
\item \textbf{（イ）集まる}\par\nobreak\noindent （ア）の矢印を全部逆向きにした流れです。入る一方なので「\mbox{出－入}」は負。どこも吸いこみです。
\item \textbf{（ウ）そろった流れ}\par\nobreak\noindent 左の辺から入った水が、同じだけ右の辺から出ます。上と下の辺は、水が辺にそって流れるだけで、通りぬけません。だから「\mbox{出－入}」は0です。
\end{enumerate}

\noindent 残りの3つの場も、同じように見ていきます。

\begin{bkfigure}
\begin{tabular}{@{}>{\centering\arraybackslash}p{40mm}@{\hspace{3mm}}>{\centering\arraybackslash}p{40mm}@{\hspace{3mm}}>{\centering\arraybackslash}p{40mm}@{}}
\kviiijpanel{-0.32*\y}{0.32*\x}{\draw[kviiij box] (0.5,0.5) rectangle (1.5,1.5);} &
\kviiijpanel{0.32*\y}{0}{\draw[kviiij box] (-0.5,0.5) rectangle (0.5,1.5);} &
\kviiijpanel{0.16*(\x+3)}{0}{\draw[kviiij box] (-0.5,-0.5) rectangle (0.5,0.5);} \\[1mm]
{\sffamily\bfseries\small（エ）うず $(-y,\ x)$} & {\sffamily\bfseries\small（オ）ずれ $(y,\ 0)$} & {\sffamily\bfseries\small（カ）側溝 $(x,\ 0)$} \\
{\sffamily\small 出＝入：発散は0} & {\sffamily\small 出＝入：発散は0} & {\sffamily\small 出＞入：発散は正} \\
\end{tabular}
\bkfigcaption{主役の場（エ）〜（カ）。（カ）は、右へ行くほど速い流れの一部を、矢印がみな右向きになるようにずらしてかいた。}
\label{bk:8-9:fig-six2}
\end{bkfigure}

\begin{enumerate}[label={\bkcircled[\bkTT]{\small\arabic*}}, leftmargin=2.4\zw, itemsep=3pt, topsep=2pt, start=4]
\item \textbf{（エ）うず}\par\nobreak\noindent 水はぐるぐる回っています。それでも、マスに入った分は、反対側の辺からちょうど同じだけ出ていきます。回るようすは、8-10 の「回転 rot」で調べます。\textbf{回っていても、発散は0}です。
\item \textbf{（オ）ずれ}\par\nobreak\noindent 上へ行くほど右向きに速い流れです。左の辺から入る量と右の辺から出る量は、同じ高さでは同じ速さなので、等しくなります。上と下の辺は、水が辺にそって流れるだけです。発散は0です。
\item \textbf{（カ）側溝}\par\nobreak\noindent 矢印はみな右向きです。ところが、右へ行くほど長い（速い）ので、右の辺から出る量は、左の辺から入る量より多くなります。発散は正です。
\end{enumerate}

\kviiijhead{向きがそろっていても、わき出しがある}
（カ）は、雨の日の道路わきのみぞ（側溝）に似ています。雨水が、みぞ全体に少しずつ流れこみます。そのため、下流ほど水が多く、速く流れます。

\begin{bkfigure}
\begin{tikzpicture}[x=1mm, y=1mm]
  \fill[bkBlueBg] (0,0) rectangle (132,20);
  \draw[bkGray, line width=1pt] (0,0) -- (132,0) (0,20) -- (132,20);
  \foreach \x/\l in {6/5,32/9,58/13,84/17,108/21} {\foreach \y in {5,15} {\draw[bk arrow, bkTheme, line width=1.3pt] (\x,\y) -- ++(\l,0);}}
  \foreach \x in {10,30,50,70,90,110} {\draw[bkBlue!60!black, line width=0.8pt] (\x,31) -- (\x-2,25);}
  \draw[kviiij box] (43,1) rectangle (56,19);
  \node[font=\sffamily\small, anchor=north, text=bkInk] at (49.5,-1.8) {左の辺から入る量 ＜ 右の辺から出る量};
  \node[font=\sffamily\small, text=bkInk, anchor=west] at (114,29.5) {雨};
\end{tikzpicture}
\bkfigcaption{側溝の流れ。点線のマスでは、右の辺から出る量のほうが多い。雨水が、マスの中で流れに加わっているから。}
\label{bk:8-9:fig-sokko}
\end{bkfigure}

矢印はみな右向きで、広がっても集まってもいません。それでも、どのマスでも出る量が入る量より多く、発散は正です。雨水が、みぞのどこでも少しずつ「わき出し」になっているのです。

\begin{taisetsu}
発散は、矢印の\textbf{向き}ではなく、\textbf{出る量と入る量の差}で決まる。\\
回っていても0（エ）。向きがそろっていても正のことがある（カ）。
\end{taisetsu}

%% ---------------- 手でわかる ----------------
\bkstage{te}{手でわかる}{マスの4つの辺で数える}{\bklevel{4}}

\bksubsection{3}{速さ×辺の長さで、通る量を数える}

\label{bk:8-9:masu}こんどは、数で計算します。1辺 $h=\qty{1}{m}$ のマスの4つの辺で、辺をまっすぐ通りぬける水の速さがわかっているとします。

\begin{bkfigure}
\begin{minipage}[c]{0.45\linewidth}\centering
\begin{tikzpicture}[x=1mm, y=1mm]
  \fill[bkOrange!8] (0,0) rectangle (28,28);
  \draw[kviiij box] (0,0) rectangle (28,28);
  \draw[bk arrow, bkTheme, line width=1.5pt] (29.5,14) -- (37.5,14);
  \node[font=\sffamily\small, text=bkInk, anchor=west] at (38.5,14) {3};
  \draw[bk arrow, bkTheme, line width=1.5pt] (-9.5,14) -- (-1.5,14);
  \node[font=\sffamily\small, text=bkInk, anchor=east] at (-10.5,14) {1};
  \draw[bk arrow, bkTheme, line width=1.5pt] (14,29.5) -- (14,40);
  \node[font=\sffamily\small, text=bkInk, anchor=south] at (14,41) {2};
  \draw[bk arrow, bkTheme, line width=1.5pt] (14,-12) -- (14,-1.5);
  \node[font=\sffamily\small, text=bkInk, anchor=north] at (14,-13) {2};
  \node[font=\sffamily\footnotesize, text=bkInk, align=center] at (14,14) {$h=\qty{1}{m}$\\速さは m/s};
\end{tikzpicture}
\end{minipage}\hfill
\begin{minipage}[c]{0.52\linewidth}
\raggedright
右の辺では、水が \qty{3}{m/s} で外へ出ています。左の辺では \qty{1}{m/s} で中へ入っています。上の辺では \qty{2}{m/s} で出て、下の辺では \qty{2}{m/s} で入っています。
\end{minipage}
\bkfigcaption{1辺\qty{1}{m}のマス。数は、辺をまっすぐ通りぬける水の速さ（m/s）。}
\label{bk:8-9:fig-masu}
\end{bkfigure}

1秒間に1つの辺を通る水の量は、「\textbf{速さ×辺の長さ}」です。\qty{3}{m/s} なら、水は1秒間に\qty{3}{m}進みます。はば\qty{1}{m}の辺を通った水は、たて\qty{3}{m}・横\qty{1}{m}の長方形の分です。8-8 の流束（速さ×面積）と同じ考え方です。平面の流れなので、量は面積（\unit{m^2}）で数えます（深さ\qty{1}{m}ぶんの水の体積と考えてもかまいません）。

出る量をプラス、入る量をマイナスにして、4つの辺を合計します。
\begin{bkhenkei}[0.6]
(+3\times1)+(-1\times1)+(+2\times1)+(-2\times1) & 出る量はプラス、入る量はマイナス \\
=3-1+2-2=2 & 1秒間の「出－入」は \qty{2}{m^2} \\
\end{bkhenkei}
マスの面積は $1\times1=\qty{1}{m^2}$ なので、1平方メートルあたりにすると、
\kviiijline{発散 $=\dfrac{\qty{2}{m^2/s}}{\qty{1}{m^2}}=2$\,\kviiijps}
正なので、このマスはわき出しです。単位の「\kviiijps」は「毎秒」と読みます。「1秒間に、1平方メートルあたり\qty{2}{m^2}の水が生まれている」という意味です。\unit{m^2} と \unit{m^2} が約分されて、\kviiijps だけが残りました。

\bksubsection{4}{右向き・上向きを正にして書く}

\label{bk:8-9:fugou}表計算でたくさんのマスを計算するには、「出る・入る」をいちいち決めるのは大変です。そこで、8-5 と同じように、$x$ 方向（右向き）の速さを $v_x$、$y$ 方向（上向き）の速さを $v_y$ と書きます。\textbf{左向きなら $v_x$ はマイナス、下向きなら $v_y$ はマイナス}です。すると、4つの辺はこうなります。

\begin{bkcenter}
\begin{bktable}{colspec={Q[l,m,wd=14mm] Q[c,m,wd=18mm] Q[l,m,wd=38mm] Q[l,m,wd=34mm]}, rowsep=1.5pt, row{1}={font=\sffamily\small}}
辺 & 外へ向かう向き & $v_x$・$v_y$ が正のとき & 外へ出る量（\mbox{出－入}への足し方） \\
右の辺 & 右 & 右へ流れる＝出る & $+(\text{右の }v_x)\times h$ \\
左の辺 & 左 & 右へ流れる＝\textbf{入る} & $-(\text{左の }v_x)\times h$ \\
上の辺 & 上 & 上へ流れる＝出る & $+(\text{上の }v_y)\times h$ \\
下の辺 & 下 & 上へ流れる＝\textbf{入る} & $-(\text{下の }v_y)\times h$ \\
\end{bktable}
\bktabcaption{左と下の辺は、外へ向かう向きが、$v_x$・$v_y$ の正の向きと反対。だから、ひく。}
\label{bk:8-9:tab-fugou}
\end{bkcenter}

左の辺で水が左向きに流れていれば、$v_x$ はマイナスです。「マイナスの数をひく」ので、たすことになり、ちゃんと「出る」の分になります。符号は、自動でそろうのです。4つをたして、面積 $h\times h$ でわると、
\kviiijframe[124mm]{発散 $\approx\dfrac{(\text{右の }v_x)-(\text{左の }v_x)}{h}+\dfrac{(\text{上の }v_y)-(\text{下の }v_y)}{h}$}
（$h\times h$ でわるとき、辺の長さ $h$ と約分されて、$h$ が1つ残ります。）「$\approx$」は「ほぼ等しい」です。マスの辺の上で速さが変わるとき、辺のまん中の速さを使うと、よい近さになります（⑤で、等号になるわけを見ます）。

図\ref{bk:8-9:fig-masu}のマスでは、$\dfrac{3-1}{1}+\dfrac{(-2)-(-2)}{1}$ ではありません。下の辺の水は上向きに流れて入っているので、下の $v_y$ は $+2$ です。だから $\dfrac{3-1}{1}+\dfrac{2-2}{1}=2$ で、前と同じ答えになります。

\bksubsection{5}{表計算で、池の口をさがす}

\label{bk:8-9:ike}いよいよ、\bkhatenaref の池です。池を1辺\qty{1}{m}のマス 6×4 に分けます。マスの辺を1秒間に通る水の量（辺にそった速さの平均×辺の長さ）を、流れの計算から求めて、表計算のシートに入れました。8-1 と同じように、シートの行は北（上）から、列は西（左）から並べてあります。

\textbf{たての辺を右向きに通る量}（$v_x$ の量。右向きが正）は、7本の線 $x=0,\ 1,\ \ldots,\ 6$ の上の、4つの段の辺です。

\begin{bkcenter}
\bksheet{7}{12}{%
1 & $x=0$ & $x=1$ & $x=2$ & $x=3$ & $x=4$ & $x=5$ & $x=6$ \\
2 & $-0.18$ & $-0.16$ & 0.45 & 0.45 & 0.30 & 0.09 & $-0.02$ \\
3 & $-0.28$ & $-0.83$ & 1.22 & 0.70 & 0.55 & $-0.11$ & $-0.15$ \\
4 & $-0.15$ & $-0.11$ & 0.55 & 0.70 & 1.22 & $-0.83$ & $-0.28$ \\
5 & $-0.02$ & 0.09 & 0.30 & 0.45 & 0.45 & $-0.16$ & $-0.18$ \\}
\bktabcaption{2行めは $y$ が3〜4の段、3行めは2〜3、4行めは1〜2、5行めは0〜1の段（単位は \unit{m^2/s}。小数第2位までに丸めて見せた）。}
\end{bkcenter}

\textbf{横の辺を上向きに通る量}（$v_y$ の量。上向きが正）は、5本の線 $y=4,\ 3,\ 2,\ 1,\ 0$ の上の、6つの列の辺です。

\begin{bkcenter}
\bksheet{6}{12}{%
6 & 0〜1 & 1〜2 & 2〜3 & 3〜4 & 4〜5 & 5〜6 \\
7 & 0.22 & 0.30 & 0.14 & $-0.06$ & $-0.17$ & $-0.17$ \\
8 & 0.24 & 0.91 & 0.14 & $-0.22$ & $-0.37$ & $-0.28$ \\
9 & $-0.32$ & $-1.04$ & $-0.37$ & $-0.37$ & $-1.04$ & $-0.32$ \\
10 & $-0.28$ & $-0.37$ & $-0.22$ & 0.14 & 0.91 & 0.24 \\
11 & $-0.17$ & $-0.17$ & $-0.06$ & 0.14 & 0.30 & 0.22 \\}
\bktabcaption{7行めは $y=4$ の線、8行めは $y=3$、9行めは $y=2$、10行めは $y=1$、11行めは $y=0$ の線（6行めの「0〜1」などは $x$ の範囲。単位は \unit{m^2/s}）。}
\end{bkcenter}

左上のマス（$x$ が0〜1、$y$ が3〜4）を考えます。右の辺は B2、左の辺は A2、上の辺は A7、下の辺は A8 です。表\ref{bk:8-9:tab-fugou}のとおり「右－左＋上－下」なので、13行めの A13 に
\kviiijline{\texttt{=(B2-A2)+(A7-A8)}}
と入れます（マスの面積は\qty{1}{m^2}なので、わる計算は省きました）。この式を、右へ F13 まで、下へ16行めまでコピーします。右へコピーすると、B と A が C と B に1つずつずれて、となりのマスの計算になります。下へコピーすると、1つ下の段のマスの計算になります。

結果を、池と同じ並びの「\textbf{発散の地図}」にしました。

\begin{bkfigure}
\begin{tikzpicture}[x=13.5mm, y=13.5mm]
  \kviiijpondmap
  \kviiijpondgrid
  \foreach \i in {0,...,6} {\node[bk tick, below=1mm] at (\i,0) {\i};}
  \foreach \j in {0,...,4} {\node[bk tick, left=1mm] at (0,\j) {\j};}
  \node[bk label, anchor=north west] at (6.15,-0.05) {$x$（m）};
  \node[bk label, anchor=south west] at (0.05,4.08) {$y$（m）};
\end{tikzpicture}
\bkfigcaption{各マスの発散（単位 \kviiijps、小数第1位まで）。オレンジのマスは正（わき出し）、青のマスは負（吸いこみ）。}
\label{bk:8-9:fig-map}
\end{bkfigure}

0でないマスは2つだけです。$x$ が1〜2、$y$ が2〜3 のマスは $+4.0$ で、ここに\textbf{わき出す口}があります。$x$ が4〜5、$y$ が1〜2 のマスは $-4.0$ で、ここに\textbf{吸いこむ口}があります。これが\bkhatenaref の答えです。

わき出す口のマスを、たしかめ算してみましょう。右 C3、左 B3、上 B8、下 B9 です。
\begin{bkhenkei}[0.48]
(1.22-(-0.83))+(0.91-(-1.04)) & 「右－左＋上－下」 \\
=1.22+0.83+0.91+1.04 & マイナスをひく＝たす \\
=4.00 & 1秒に \qty{4}{m^2} わき出す \\
\end{bkhenkei}
左の辺（$-0.83$）と下の辺（$-1.04$）は、マイナスです。水が左向き・下向き、つまりマスの外へ流れているからです。4つの辺のどれからも、水は外へ出ています。

\kviiijhead{全部のマスをたすと？}
池の24マスの発散を全部たすと、$4.0+(-4.0)=0$ です。いっぽう、池の外まわり（いちばん外の辺）を通って出る量を全部たしても、0 になります（A列とG列、7行めと11行めの数で確かめられます）。となりどうしのマスでは、共通の辺の水が「片方から出て、もう片方に入る」ので、たすと打ち消し合うからです。

\begin{bkfigure}
\begin{tikzpicture}[x=1mm, y=1mm]
  \fill[bkOrange!20] (0,0) rectangle (26,26);
  \fill[bkBlueBg] (26,0) rectangle (52,26);
  \draw[bkGray, line width=0.8pt] (0,0) rectangle (52,26) (26,0) -- (26,26);
  \draw[kviiij big] (20,13) -- (32,13);
  \node[font=\sffamily\small, text=bkInk, anchor=south west] at (27.5,14.5) {1.22};
  \node[font=\sffamily\footnotesize, text=bkInk, align=center] at (11,6) {口のマス};
  \node[font=\sffamily\footnotesize, text=bkInk, align=center] at (41,6) {右どなり};
  \node[font=\sffamily\small, text=bkInk, align=left, anchor=west] at (58,13) {口のマス：右の辺から出る $+1.22$\\右どなり：左の辺から入る $-1.22$\\2つをたすと $+1.22-1.22=0$};
\end{tikzpicture}
\bkfigcaption{となりどうしのマスの共通の辺（池の、$x=2$・$y$ が2〜3の辺）。片方の「出る」は、もう片方の「入る」。}
\label{bk:8-9:fig-tonari}
\end{bkfigure}

マスの発散の合計が、外まわりから出る量に等しい。このことは 8-11 の「ガウスの定理」で、くわしく学びます。

\kviiijhead{コンピュータで、マスの出入りを数える}
辺を通る量は、「辺を細かく分けて、速さ×小さな長さをたす」と求められます（8-8 の、データだけのときのやり方と同じ）。Python で、池の流れの式から、24マスの発散を計算してみます。プログラムの中の $1.5,\ 2.5$ と $4.5,\ 1.5$ は口の場所、$4.0$ は1秒間にわき出す量 $Q$ です。式の形（なぜ $2\pi$ でわるのか）は、8-11 で学びます。

\Needspace{18\baselineskip}
\begin{bkcode}[Python]
import numpy as np
def v(x, y):    \# 池の流れの速さ (vx, vy)
    ax, ay = x - 1.5, y - 2.5    \# わき出す口から見た場所
    bx, by = x - 4.5, y - 1.5    \# 吸いこむ口から見た場所
    a2, b2 = ax**2 + ay**2, bx**2 + by**2
    k = 4.0 / (2 * np.pi)        \# Q = 4
    return k*(ax/a2 - bx/b2), k*(ay/a2 - by/b2)
t = (np.arange(100) + 0.5) / 100    \# 辺を100こに分けた、まん中の点
for j in [3, 2, 1, 0]:              \# 北（上）の段から
    row = ""
    for i in range(6):
        d = (v(i+1, j+t)[0].mean() - v(i, j+t)[0].mean()
             + v(i+t, j+1)[1].mean() - v(i+t, j)[1].mean())
        row += f"\{round(d, 1) + 0.0:6.1f\}"
    print(row)
\end{bkcode}
\begin{bkcodeout}[実行結果]
   0.0   0.0   0.0   0.0   0.0   0.0
   0.0   4.0   0.0   0.0   0.0   0.0
   0.0   0.0   0.0   0.0  -4.0   0.0
   0.0   0.0   0.0   0.0   0.0   0.0
\end{bkcodeout}
\vspace{4pt}
\noindent \texttt{t} は、0から1までを100こに分けた小さな区間の、まん中の点です。\texttt{v(i+1, j+t)[0]} は、マスの右の辺（$x=i+1$）の上の100か所の $v_x$ で、\texttt{.mean()} はその平均です。平均に辺の長さ\qty{1}{m}をかけると、辺を通る量になります（1をかけても変わらないので、かけていません）。\texttt{d} が「右－左＋上－下」です。\texttt{+ 0.0} は、$-0.0$ と表示されないようにするためのくふうです。表計算と同じ地図が出てきました。

\bkkokomade[小さなマスを置いて「出－入」を数え、表計算でわき出しと吸いこみの場所を見つけられれば十分です。この先は、マスを小さくしていって、式で発散を求めます。]

%% ---------------- 式でわかる ----------------
\Needspace{9\baselineskip}\bkstage{shiki}{式でわかる}{マスを小さくしていく}{\bklevel{5}}

\bksubsection{6}{マスを小さくすると、偏微分の和になる}

\label{bk:8-9:shiki}マスを大きくとると、「マスの中のどこに口があるか」まではわかりません。マスを小さくしていくほど、場所ごとの発散がくわしくわかります。そこで、マスをどこまでも小さくしたときに、発散の式がどうなるかを調べます。

マスの中心を $(x,\ y)$、1辺を $h$ とします。右の辺は中心から右へ $\frac h2$、左の辺は左へ $\frac h2$ の所にあります。辺のまん中の速さを使って、④の式を書き直します。
\begin{bkhenkei}[0.6]
\frac{(\text{右の }v_x)-(\text{左の }v_x)}{h} & ④の式の1つめの項 \\
=\frac{v_x\bigl(x+\frac h2,\ y\bigr)-v_x\bigl(x-\frac h2,\ y\bigr)}{h} & 右の辺のまん中は $\bigl(x+\frac h2,\ y\bigr)$、左は $\bigl(x-\frac h2,\ y\bigr)$ \\
\end{bkhenkei}
分子の2つの点は、$y$ が同じで、$x$ だけがちがいます。$x$ のちがいは、ちょうど $h$ です。だから、この分数は「\textbf{$y$ を止めて、$x$ だけを $h$ 動かしたときの、$v_x$ の変化の割合}」です。$h$ を0に近づけると、8-2 で学んだ偏微分 $\dfrac{\partial v_x}{\partial x}$ になります。

2つめの項も同じです。上の辺と下の辺のまん中は $\bigl(x,\ y+\frac h2\bigr)$ と $\bigl(x,\ y-\frac h2\bigr)$ で、$x$ が同じで $y$ だけがちがいます。$h$ を0に近づけると $\dfrac{\partial v_y}{\partial y}$ になります。

\Needspace{18\baselineskip}\kviiijhead{数で確かめる：$h$ を小さくすると、近づく}
$v_x=x^3$ の流れで、$x=1$ の所の1つめの項を、$h$ を小さくしながら計算します。
\begin{bkcenter}
\begin{bktable}{colspec={Q[c,m,wd=12mm] Q[c,m,wd=60mm] Q[c,m,wd=24mm]}, rowsep=3pt, row{1}={font=\sffamily\small}}
$h$ & $\dfrac{(1+\frac h2)^3-(1-\frac h2)^3}{h}$ & 計算の答え \\
1 & $\dfrac{1.5^3-0.5^3}{1}=\dfrac{3.375-0.125}{1}$ & 3.25 \\
0.5 & $\dfrac{1.25^3-0.75^3}{0.5}$ & 3.0625 \\
0.1 & $\dfrac{1.05^3-0.95^3}{0.1}$ & 3.0025 \\
0.01 & $\dfrac{1.005^3-0.995^3}{0.01}$ & 3.000025 \\
\end{bktable}
\bktabcaption{$h$ を小さくすると、3 に近づく。}
\end{bkcenter}
いっぽう、$v_x=x^3$ を $x$ で微分すると $3x^2$、$x=1$ では $3\times1^2=3$ です。表の数は、ちょうどこの3に近づいています。（式で計算すると、表の数は $3+\frac{h^2}{4}$ になっています。$h=0.1$ なら $3+0.0025$。$h$ が小さいほど、おまけの $\frac{h^2}{4}$ は小さくなります。）

\begin{taisetsu}
平面の流れ $\vec v=(v_x,\ v_y)$ の発散は
\kviiijline{\large $\mathrm{div}\,\vec v=\dfrac{\partial v_x}{\partial x}+\dfrac{\partial v_y}{\partial y}$}
\textbf{$v_x$ は $x$ で、$v_y$ は $y$ で}偏微分して、たす。答えは、場所ごとに1つの数。
\end{taisetsu}

「$v_x$ は $x$ で」の組み合わせには、わけがあります。$x$ 方向（左右）の出入りは、左右の辺の $v_x$ のちがいで決まりました。左右の辺は、$x$ だけがちがう2つの場所です。だから、$v_x$ を $x$ で微分します。$v_x$ を $y$ で微分した量（上下で $v_x$ がどう変わるか）は、出入りには関係しません。それは、（オ）のずれのように、辺にそって流れる水の速さのちがいだからです。

\Needspace{12\baselineskip}
\bksubsection{7}{主役の場で、目と式をくらべる}

\label{bk:8-9:hyo}図\ref{bk:8-9:fig-six}・図\ref{bk:8-9:fig-six2}の6つの場を、式で計算します。たとえば（ア）$(x,\ y)$ では、$v_x=x$、$v_y=y$ なので、$\dfrac{\partial v_x}{\partial x}=\dfrac{\partial x}{\partial x}=1$、$\dfrac{\partial v_y}{\partial y}=\dfrac{\partial y}{\partial y}=1$。たして2です。

\begin{bkcenter}
\begin{bktable}{colspec={Q[l,m,wd=27mm] Q[c,m,wd=20mm] Q[c,m,wd=13mm] Q[c,m,wd=13mm] Q[c,m,wd=13mm] Q[c,m,wd=14mm]}, colsep=4pt, rowsep=1.5pt, row{1}={font=\sffamily\small}}
流れ & $(v_x,\ v_y)$ & $\dfrac{\partial v_x}{\partial x}$ & $\dfrac{\partial v_y}{\partial y}$ & $\mathrm{div}\,\vec v$ & 目で見た結果 \\
（ア）広がる & $(x,\ y)$ & 1 & 1 & 2 & 正 \\
（イ）集まる & $(-x,\ -y)$ & $-1$ & $-1$ & $-2$ & 負 \\
（ウ）そろった流れ & $(1,\ 0)$ & 0 & 0 & 0 & 0 \\
（エ）うず & $(-y,\ x)$ & 0 & 0 & 0 & 0 \\
（オ）ずれ & $(y,\ 0)$ & 0 & 0 & 0 & 0 \\
（カ）側溝 & $(x,\ 0)$ & 1 & 0 & 1 & 正 \\
\end{bktable}
\bktabcaption{目で見た結果と、式で計算した結果が一致する。}
\label{bk:8-9:tab-six}
\end{bkcenter}

（エ）の $v_x=-y$ には $x$ が入っていないので、$x$ で偏微分すると0です（$y$ を止めると、$-y$ はただの数）。$v_y=x$ を $y$ で偏微分しても0です。（オ）も同じです。

（ア）では、マスの大きさによらず、右の辺と左の辺の $v_x$ の差がいつも $h$ です（$\bigl(x+\frac h2\bigr)-\bigl(x-\frac h2\bigr)=h$）。だから、④の近い式でも $\frac hh+\frac hh=2$ となり、式の答えとぴったり同じです。$v_x$・$v_y$ が1次式なら、マスが大きくても④の式は「$\approx$」でなく「$=$」になります。

\bksubsection{8}{空間（3次元）の流れ}

\label{bk:8-9:sanjigen}空間では、マスのかわりに1辺 $h$ の立方体の箱を置きます。面は6つ。向かい合う面が3組あります。

\begin{bkfigure}
\begin{tikzpicture}[x=1mm, y=1mm]
  % 立方体（斜めから見た図）
  \def\kviiijd{9}
  \fill[bkOrange!25] (30,0) -- (30,30) -- (30+\kviiijd,30+\kviiijd*0.6) -- (30+\kviiijd,\kviiijd*0.6) -- cycle;
  \fill[bkOrange!25] (0,0) rectangle (30,30);
  \draw[bkInk, line width=0.8pt] (0,0) rectangle (30,30);
  \draw[bkInk, line width=0.8pt] (0,30) -- (\kviiijd,30+\kviiijd*0.6) -- (30+\kviiijd,30+\kviiijd*0.6) -- (30,30);
  \draw[bkInk, line width=0.8pt] (30+\kviiijd,30+\kviiijd*0.6) -- (30+\kviiijd,\kviiijd*0.6) -- (30,0);
  \draw[bkInk!60, line width=0.6pt, dash pattern=on 2pt off 1.5pt] (0,0) -- (\kviiijd,\kviiijd*0.6) -- (30+\kviiijd,\kviiijd*0.6) (\kviiijd,\kviiijd*0.6) -- (\kviiijd,30+\kviiijd*0.6);
  % x の向きの矢印（左の面から入り、右の面から出る）
  \draw[kviiij big] (-16,15) -- (-2,15);
  \draw[kviiij big] (41,17.7) -- (54,17.7);
  \node[font=\sffamily\small, text=bkInk, anchor=south] at (-9,16.5) {入る};
  \node[font=\sffamily\small, text=bkInk, anchor=south] at (47.5,19.2) {出る};
  % 軸
  \draw[bk axis] (64,4) -- (80,4) node[right, bk label] {$x$};
  \draw[bk axis] (64,4) -- (64,20) node[above, bk label] {$z$};
  \draw[bk axis] (64,4) -- (73,9.4) node[above right, bk label, inner sep=1pt] {$y$};
  \node[font=\sffamily\small, text=bkInk, align=left, anchor=west] at (85,15) {向かい合う面の組\\① 左右の面：$v_x$ の差\\② 前後の面：$v_y$ の差\\③ 上下の面：$v_z$ の差};
\end{tikzpicture}
\bkfigcaption{1辺 $h$ の立方体の箱。オレンジの左右の面では、$v_x$ が出入りを決める。前後・上下の面も同じように数える。}
\label{bk:8-9:fig-cube}
\end{bkfigure}

1つの面の面積は $h^2$ です。左右の面の組では、「出－入」は次のようになります。
\kviiijline{$\bigl\{(\text{右の }v_x)-(\text{左の }v_x)\bigr\}\times h^2$}
前後・上下の組も同じように数えて3つをたし、箱の体積 $h^3$ でわります。$h^2\div h^3=\frac1h$ なので、
\begin{align*}
\text{発散}&\approx\frac{(\text{右の }v_x)-(\text{左の }v_x)}{h}+\frac{(\text{奥の }v_y)-(\text{手前の }v_y)}{h}\\
&\quad+\frac{(\text{上の }v_z)-(\text{下の }v_z)}{h}
\end{align*}
$h\to0$ で、平面と同じく偏微分の和になります。
\begin{taisetsu}
空間の流れ $\vec v=(v_x,\ v_y,\ v_z)$ の発散は
\kviiijline{\large $\mathrm{div}\,\vec v=\dfrac{\partial v_x}{\partial x}+\dfrac{\partial v_y}{\partial y}+\dfrac{\partial v_z}{\partial z}$}
平面の式に、$z$ の項が1つふえるだけ。
\end{taisetsu}
物理の本（力-13）の「$\mathrm{div}\,\vec v=\dfrac{\partial v_x}{\partial x}+\dfrac{\partial v_y}{\partial y}+\dfrac{\partial v_z}{\partial z}=0$」は、この式です。

\Needspace{14\baselineskip}
\bksubsection{9}{ナブラで書く：$\nabla\cdot\vec v$}

\label{bk:8-9:nabla}8-4 で、ナブラ $\nabla$ を、偏微分を並べた「微分の矢印」として使いました。
\kviiijline{$\nabla=\Bigl(\dfrac{\partial}{\partial x},\ \dfrac{\partial}{\partial y},\ \dfrac{\partial}{\partial z}\Bigr)$}
この $\nabla$ と $\vec v=(v_x,\ v_y,\ v_z)$ の内積（4-12）を、ふつうのベクトルと同じ形で計算してみます。内積は「成分どうしをかけて、たす」でした。ただし、$\nabla$ の成分は数ではなく「〜で偏微分する」という\textbf{操作}なので、「かける」を「右の量を偏微分する」と読みかえます。
\begin{bkhenkei}[0.52]
\nabla\cdot\vec v=\frac{\partial}{\partial x}v_x+\frac{\partial}{\partial y}v_y+\frac{\partial}{\partial z}v_z & 内積の形：1つめどうし、2つめどうし、3つめどうし \\
=\frac{\partial v_x}{\partial x}+\frac{\partial v_y}{\partial y}+\frac{\partial v_z}{\partial z} & 「かける」を「偏微分する」と読んだ \\
=\mathrm{div}\,\vec v & 発散の式と同じ \\
\end{bkhenkei}
内積の形にすると、「1つめどうし」なので $v_x$ は $\frac{\partial}{\partial x}$ と組になります。成分と文字の組み合わせを、自然にまちがえない書き方です。この本では $\mathrm{div}\,\vec v$ を主に使い、すぐ横に $\nabla\cdot\vec v$ も並べて書きます。専門書では、どちらの書き方も使われます。

\kviiijhead{grad と div：矢印と数の行き来}
8-4 の勾配 $\mathrm{grad}\,f=\nabla f$ と、この章の発散 $\mathrm{div}\,\vec v=\nabla\cdot\vec v$ は、記号がにています。でも、入れる物と出てくる物が反対です。
\begin{bkcenter}
\begin{bktable}{colspec={Q[l,m,wd=26mm] Q[c,m,wd=27mm] Q[c,m,wd=22mm] Q[l,m,wd=34mm]}, rowsep=1.5pt, row{1}={font=\sffamily\small}}
記号 & 入れる物 & 出てくる物 & たとえ \\
$\mathrm{grad}\,f=\nabla f$ & 数の場（高さ・温度） & 矢印の場 & いちばん急な上り坂の矢印 \\
$\mathrm{div}\,\vec v=\nabla\cdot\vec v$ & 矢印の場（流れ） & 数の場 & わき出しの強さ \\
\end{bktable}
\bktabcaption{grad は数から矢印を、div は矢印から数を作る。$\nabla$ のうしろに「・」があれば内積なので、答えは数。}
\end{bkcenter}

\bkkigocards{\bkkigocard{$\mathrm{div}$}{ダイバージェンス}{発散。小さな箱から出ていく量（出－入）を、箱の大きさでわったもの。答えは数}{divergence}{8-9}}%
{\bkkigocard{$\nabla\cdot$}{ナブラ ドット}{$\mathrm{div}$ と同じ。$\nabla$ と矢印の場の内積の形で書く}{nabla dot（del dot）}{8-9}}

\bksubsection{10}{流束で決める：発散のせいかくな定義}

\label{bk:8-9:teigi}ここまで、箱を「マス」や「立方体」にしてきました。では、箱の形を球や円柱にしたら、ちがう答えになるのでしょうか。また、8-8 で学んだ「閉じた面の流束」を使うと、発散はどう書けるでしょうか。

8-8 では、閉じた面 $S$ から外へ出ていく量（流束）を
\[ \kviiijNt{\oiint_S\vec v\cdot\vec n\,\dd S} \]
と書きました。オレンジの点線の枠と「名前」の札は、この式が\textbf{操作の名前}で、まだ計算できる形ではないという印です（8-6・8-8 と同じ）。「閉じた面 $S$ を小さなタイルに分け、1枚ずつ『$\vec v$ の外向きの成分×タイルの面積』を求めて、全部たし合わせる」という操作の名前です。$\oiint$ の◯は「\textbf{閉じている}」という印で、球の形や面積を表すのではありません。どの面かは、下の $S$ が決めます。

この「出－入」を、箱の体積 $\Delta V$ でわって、箱を小さくしていった先が、発散です。

\begin{seikaku}
\kviiijbox{発散の定義}
点Pを囲む小さな閉じた面 $S$（中の体積 $\Delta V$）を考え、$S$ から外へ出ていく流束を $\Delta V$ でわります。$S$ を点Pに向けて小さくしていったときに近づく値を、点Pの発散といいます。
\kviiijline{$\displaystyle(\mathrm{div}\,\vec v)(\mathrm{P})=\lim_{\Delta V\to0}\frac{1}{\Delta V}\oiint_S\vec v\cdot\vec n\,\dd S$}
$\vec n$ は面に直角で外向きの、長さ1の矢印です（8-8）。

\kviiijbox{実際の操作}
この式は「名前」なので、そのままでは計算できません。実際には、8-8 の場合分けの表から、やり方を選びます。$S$ を立方体にとると、場合分けの表8-8の \kviiijcase{G}（閉じた面が平らな面の組み合わせ：面ごとに求めてたす）です。左右・前後・上下の6つの面で $\vec v\cdot\vec n$ は $\pm v_x$、$\pm v_y$、$\pm v_z$ になり、⑧の計算そのものになります。だから、この定義から $\dfrac{\partial v_x}{\partial x}+\dfrac{\partial v_y}{\partial y}+\dfrac{\partial v_z}{\partial z}$ が出てきます（例題6で、面ごとに計算します）。

\kviiijbox{箱の形によらない}
$v_x$・$v_y$・$v_z$ がなめらかに変わる（偏微分できて、その偏微分も連続な）場なら、$S$ を球や円柱など、どんな形にしても同じ値に近づきます。この定義は座標の軸の向きにもよらないので、物理の本では、この形で発散を考えることがよくあります。
\end{seikaku}

%% ---------------- 点のわき出し・物理の発散 ----------------
\bksubsection{11}{点のわき出し：口の外では発散0}

\label{bk:8-9:ten}池の発散の地図では、口のある2マスだけが0でなく、ほかのマスは全部0でした。では、口のあるマスを、さらに小さなマスに分けたら、どうなるでしょう。

\begin{bkfigure}
\begin{minipage}[c]{0.44\linewidth}\centering
\begin{tikzpicture}[x=1mm, y=1mm]
  \fill[bkOrange!90!black] (0,0) circle (1.3);
  \foreach \a in {0,30,...,330} {\draw[bk arrow, bkTheme, line width=1pt] (\a:4) -- (\a:13);
    \draw[bk arrow, bkTheme, line width=1pt] (\a:19.5) -- (\a:24);}
  \draw[kviiij box] (15:15.5) arc (15:45:15.5) -- (45:27) arc (45:15:27) -- cycle;
  \node[font=\sffamily\footnotesize, text=bkInk, anchor=north] at (0,-27) {口から $r$ の所で、速さは $\dfrac{Q}{2\pi r}$};
\end{tikzpicture}
\end{minipage}\hfill
\begin{minipage}[c]{0.53\linewidth}
\raggedright
1つの口から1秒間に $Q$ の水がわき出すと、水は四方へ同じように広がります。口から $r$ はなれた円（周の長さ $2\pi r$）を、1秒間に $Q$ の水が通りぬけます。だから、そこでの速さは $\dfrac{Q}{2\pi r}$ です。遠くほど遅く、矢印は短くなります。
\end{minipage}
\bkfigcaption{1つの口からのわき出し。点線のように口をふくまない箱では、内側の弧から入る量と、外側の弧から出る量が等しい。}
\label{bk:8-9:fig-ten}
\end{bkfigure}

点線の箱を見てください。外側の弧は、内側の弧より長いです。でも、外側では矢印が短い（遅い）ので、「速さ×長さ」は同じになります。$\dfrac{Q}{2\pi r}\times2\pi r=Q$ で、半径によらないからです。口をふくまない箱では、「\mbox{出－入}」は0です。

式でも確かめます。口を原点に置くと、速さの矢印は $\vec v=\dfrac{Q}{2\pi}\dfrac{(x,\ y)}{x^2+y^2}$ です（向きは原点から外向き、長さは $\dfrac{Q}{2\pi r}$）。$\dfrac{Q}{2\pi}$ はただの数なので、$\dfrac{(x,\ y)}{x^2+y^2}$ の発散を計算します。商の微分（5-6 のもっと知りたい：$\bigl(\frac fg\bigr)'=\frac{f'g-fg'}{g^2}$）を使います。
\begin{bkhenkei}[0.56]
\frac{\partial}{\partial x}\Bigl(\frac{x}{x^2+y^2}\Bigr)=\frac{1\times(x^2+y^2)-x\times2x}{(x^2+y^2)^2} & $y$ を止めて $x$ で微分（分母を $x$ で微分すると $2x$） \\
=\frac{y^2-x^2}{(x^2+y^2)^2} & 分子を整理した \\
\frac{\partial}{\partial y}\Bigl(\frac{y}{x^2+y^2}\Bigr)=\frac{x^2-y^2}{(x^2+y^2)^2} & 同じように、$x$ を止めて $y$ で微分 \\
\mathrm{div}\,\vec v=\frac{Q}{2\pi}\cdot\frac{(y^2-x^2)+(x^2-y^2)}{(x^2+y^2)^2}=0 & たすと、分子が0 \\
\end{bkhenkei}
原点以外のどこでも、発散は0です。原点では $x^2+y^2=0$ で、分母が0になるので計算できません。

つまり、池の「4」は、マスの中のどこかに広がっているのではなく、\textbf{口の1点だけに集まっている}のです。口をふくむマスを、どんなに小さくしても、出る量は $Q=4$ のままです。面積でわると、マスが小さくなるほど、発散はどこまでも大きくなります。このような1点だけのわき出しを、\textbf{点のわき出し}といいます。点のわき出しは、8-11 で「口のまわりを小さな円でくりぬく」方法で調べます。

\kviiijhead{空間の点のわき出し：点電荷の電場}
空間でも同じです。プラスの電気をもつ小さな粒（点電荷）のまわりの電場は、原点から外向きで、大きさは距離 $r$ の2乗に反比例します（8-5）。向きと大きさをまとめると、$\dfrac{(x,\ y,\ z)}{r^3}$ に比例します（$r=\sqrt{x^2+y^2+z^2}$。長さ $r$ の矢印 $(x,\ y,\ z)$ を $r^3$ でわると、長さは $\frac1{r^2}$）。これも、原点以外では発散が0です。
\begin{bkhenkei}[0.62]
\frac{\partial}{\partial x}\Bigl(\frac{x}{r^3}\Bigr)=\frac{1}{r^3}-\frac{3x^2}{r^5} & 積の微分（5-6）。下の説明 \\
\frac{\partial}{\partial y}\Bigl(\frac{y}{r^3}\Bigr)=\frac{1}{r^3}-\frac{3y^2}{r^5} & 同じ形 \\
\frac{\partial}{\partial z}\Bigl(\frac{z}{r^3}\Bigr)=\frac{1}{r^3}-\frac{3z^2}{r^5} & 同じ形 \\
\text{3つの和}=\frac{3}{r^3}-\frac{3(x^2+y^2+z^2)}{r^5} & 3行をたした \\
=\frac{3}{r^3}-\frac{3r^2}{r^5}=0 & $x^2+y^2+z^2=r^2$ \\
\end{bkhenkei}
1行めの計算を、ていねいに見ます。$\dfrac{x}{r^3}=x\times r^{-3}$ と見て、積の微分を使います。まず、$r=\sqrt{x^2+y^2+z^2}$ を $x$ で微分すると、合成関数の微分（5-6）で $\dfrac{1}{2\sqrt{x^2+y^2+z^2}}\times2x=\dfrac{x}{r}$ です。だから $r^{-3}$ を $x$ で微分すると $-3r^{-4}\times\dfrac{x}{r}=-\dfrac{3x}{r^5}$ です。積の微分は「$x$ の微分×$r^{-3}$ ＋ $x$×$r^{-3}$ の微分」なので、$1\times\dfrac{1}{r^3}+x\times\Bigl(-\dfrac{3x}{r^5}\Bigr)=\dfrac{1}{r^3}-\dfrac{3x^2}{r^5}$ です。

電場のわき出しは、電荷のある点だけにあります。電荷のない所では、電場の発散は0です。

\Needspace{12\baselineskip}
\bksubsection{12}{連続の式：物は勝手に生まれも消えもしない}

\label{bk:8-9:renzoku}空気のように縮む物の流れでは、場所ごとに密度（\qty{1}{m^3}あたりの質量）$\rho$（ロー）がちがい、時間とともに変わることもあります。小さな箱（体積 $\Delta V$）について、質量の出入りを考えます。

\begin{enumerate}[label={\bkcircled[\bkTT]{\arabic*}}, leftmargin=2.4\zw, itemsep=2pt, topsep=2pt]
\item 箱の中の質量は $\rho\,\Delta V$ です。1秒間にふえる量は $\dfrac{\partial\rho}{\partial t}\Delta V$ です（場所を止めて、時間だけ動かした変化なので $\partial$）。
\item 1秒間に、1平方メートルの面を通る質量は「密度×速さ」です。だから、質量の流れの矢印は $\rho\vec v$ です。箱から1秒間に出ていく質量（\mbox{出－入}）は、発散の決め方から $\mathrm{div}(\rho\vec v)\times\Delta V$ です。
\item 質量は、生まれも消えもしません。だから、箱の中でふえる量は、入ってきた量（出ていく量のマイナス）と等しくなります。
\end{enumerate}
③を式にして、$\Delta V$ でわります。
\begin{bkhenkei}[0.5]
\frac{\partial\rho}{\partial t}\Delta V=-\mathrm{div}(\rho\vec v)\,\Delta V & ふえる量＝－（出ていく量） \\
\frac{\partial\rho}{\partial t}+\mathrm{div}(\rho\vec v)=0 & $\Delta V$ でわって、右辺を左へうつした \\
\end{bkhenkei}
これを\textbf{連続の式}（れんぞくのしき）といいます。\bkidx{れんぞくのしき}{連続の式}物理の本（力-13）と同じことばで、「\textbf{その点の密度が増える速さ＋その点からのわき出し＝0}」と読みます。

水のように、ほとんど縮まない（密度がどこでも同じで変わらない）物では、$\dfrac{\partial\rho}{\partial t}=0$ で、$\rho$ は数として外に出せます。$\rho$ でわると
\kviiijframe[64mm]{$\mathrm{div}\,\vec v=0$}
縮まないことを\textbf{非圧縮}（ひあっしゅく）といいます。\bkidx{ひあっしゅく}{非圧縮}非圧縮の流れには、どこにも、わき出しも吸いこみもありません。（池の口は、池の底から水を足したり取ったりしているので、例外の点です。）


\begin{motto}[5]
\kviiijbox{密度が場所で変わるとき：$\mathrm{div}(\rho\vec v)$ をほどく}
$\rho\vec v=(\rho v_x,\ \rho v_y,\ \rho v_z)$ の発散を、積の微分（5-6）で1つずつ計算します。$\rho$ も $v_x$ も場所で変わるので、両方を微分します。
\begin{bkhenkei}[0.64]
\frac{\partial(\rho v_x)}{\partial x}=\frac{\partial\rho}{\partial x}v_x+\rho\frac{\partial v_x}{\partial x} & 積の微分。$y$・$z$ も同じ形 \\
\mathrm{div}(\rho\vec v)=\Bigl(v_x\frac{\partial\rho}{\partial x}+v_y\frac{\partial\rho}{\partial y}+v_z\frac{\partial\rho}{\partial z}\Bigr) & 3つをたして、\\
\qquad\qquad+\rho\Bigl(\frac{\partial v_x}{\partial x}+\frac{\partial v_y}{\partial y}+\frac{\partial v_z}{\partial z}\Bigr) & 2つの組に分けた \\
=\vec v\cdot\mathrm{grad}\,\rho+\rho\,\mathrm{div}\,\vec v & 1つめは内積（4-12・8-4） \\
\end{bkhenkei}
1つめの項は「密度のちがう所から、流れが物を運んでくる分」、2つめは「流れが広がって、うすまる分」です。連続の式は $\dfrac{\partial\rho}{\partial t}+\vec v\cdot\mathrm{grad}\,\rho=-\rho\,\mathrm{div}\,\vec v$ と書けます。左辺は、8-13 で学ぶ\textbf{物質微分}（水のつぶといっしょに動きながら見た変化）です。つまり「つぶといっしょに動いて見ると、流れが広がる所（$\mathrm{div}\,\vec v>0$）では密度が下がる」という意味です。
\end{motto}

\kviiijhead{縮まない流れで、速くなるわけ}
ホースの先を指でつまむと、水は勢いよく飛び出します。ホースの中の水は縮まないので、$\mathrm{div}\,\vec v=0$ です。

\begin{bkfigure}
\begin{tikzpicture}[x=1mm, y=1mm]
  \fill[bkBlueBg] (0,4) -- (52,4) -- (82,11) -- (110,11) -- (110,19) -- (82,19) -- (52,26) -- (0,26) -- cycle;
  \draw[bkGray, line width=1pt] (0,4) -- (52,4) -- (82,11) -- (110,11) (0,26) -- (52,26) -- (82,19) -- (110,19);
  \foreach \y in {9,15,21} {\draw[bk arrow, bkTheme, line width=1.2pt] (6,\y) -- (16,\y);}
  \foreach \y in {13,17} {\draw[bk arrow, bkTheme, line width=1.2pt] (90,\y) -- (106,\y);}
  \draw[bk arrow, bkTheme, line width=1.1pt] (56,22.5) -- (68,19.2);
  \draw[bk arrow, bkTheme, line width=1.1pt] (56,7.5) -- (68,10.8);
  \draw[kviiij box] (28,1.5) -- (28,28.5) -- (86,28.5) -- (86,1.5) -- cycle;
  \node[font=\sffamily\footnotesize, text=bkInk, anchor=north] at (14,-0.5) {断面積 $A_1$、速さ $v_1$};
  \node[font=\sffamily\footnotesize, text=bkInk, anchor=north] at (104,8.5) {断面積 $A_2$、速さ $v_2$};
  \node[font=\sffamily\footnotesize, text=bkInk, anchor=south] at (57,29.5) {点線の袋：入る量＝出る量};
\end{tikzpicture}
\bkfigcaption{ホースの先が細くなる所。流れは横に集まり（$\dfrac{\partial v_y}{\partial y}$ が負）、そのぶん前に速くなる（$\dfrac{\partial v_x}{\partial x}$ が正）。}
\label{bk:8-9:fig-hose}
\end{bkfigure}

先が細くなる所では、水は横から押されて、まん中へ集まります。横の向きの $\dfrac{\partial v_y}{\partial y}$（と $\dfrac{\partial v_z}{\partial z}$）は負です。発散が0なので、そのぶん $\dfrac{\partial v_x}{\partial x}$ が正になり、水は前へ進むほど速くなります。

ホース全体で見ると、点線の袋の中のどこでも発散は0なので、袋から出る量と入る量は等しくなります（池の外まわりと同じ。8-11）。管の壁からは水が出入りしないので、入口と出口だけを数えればよく、
\kviiijframe[90mm]{$v_1A_1=v_2A_2$\qquad（1秒間に入る量＝出る量）}
断面積が半分なら、速さは2倍になります。物理の本（力-13）の $A_1v_1=A_2v_2$ と同じ式です（例題8）。

\bksubsection{13}{熱の流れと電場の発散}

\label{bk:8-9:netsu}発散は、水の流れだけの道具ではありません。「何かが流れていく矢印の場」なら、どれにでも使えます。

\kviiijhead{熱の流れ $\vec q$}
8-4 で、熱の流れの矢印 $\vec q=-k\,\mathrm{grad}\,T$（\qty{1}{m^2}を1秒間に通る熱。単位は \unit{W/m^2}）を学びました。$\mathrm{div}\,\vec q$ は、\qty{1}{m^3}あたり、1秒間に外へ出ていく熱（\unit{W/m^3}）です。単位は $\unit{W/m^2}\div\unit{m}=\unit{W/m^3}$ になります（偏微分は「÷長さ」の量だから）。

\begin{itemize}[label={\color{\bkTT}●}, leftmargin=1.6\zw, itemsep=2pt, topsep=2pt]
\item 6-3 のかべ（厚さ\qty{0.2}{m}、両面が20\,℃と0\,℃）では、どこでも $q_x=\qty{160}{W/m^2}$ で一定でした。$\dfrac{\partial q_x}{\partial x}=0$ なので、発散は0。かべの中で熱は生まれず、入った熱がそのまま出ていきます。
\item 電気を流して熱を出している板（電熱線やヒーターの板）では、熱がその場で生まれます。板の中で $q_x=5000x$〔W/m²〕（$x$ は板のまん中からの距離〔m〕）なら、$\mathrm{div}\,\vec q=\qty{5000}{W/m^3}$。\qty{1}{L}（$\qty{0.001}{m^3}$）あたり\qty{5}{W}の熱が生まれ、両側へ流れ出ています。
\end{itemize}
熱が生まれる所で $\mathrm{div}\,\vec q$ は正、熱がたまって温度が上がる所では負になります。この考えで、8-13 の熱方程式を作ります。

\kviiijhead{電場 $\vec E$ と磁場 $\vec B$}
電場の発散は、その場所の電荷の密度 $\rho$（\qty{1}{m^3}あたりの電気の量。密度と同じ文字を使います）に比例します。
\kviiijline{$\mathrm{div}\,\vec E=\dfrac{\rho}{\varepsilon_0}$\qquad （$\varepsilon_0$ は自然の定数：真空の誘電率）}
プラスの電荷がある所は電場のわき出し、マイナスの電荷の所は吸いこみです。電荷のない所では0で、⑪の点電荷の計算と同じです。いっぽう、磁場の発散はどこでも0です（$\mathrm{div}\,\vec B=0$）。磁石を細かく切っても、N極だけ、S極だけの磁石はできません。磁力線にはわき出しも吸いこみもなく、いつも輪になってつながっています。2つの式は、8-14 で「マクスウェル方程式」として読みます。



\begin{motto}[5]
\kviiijbox{流れをたし合わせると、発散もたし合わせになる}
池の流れは、口 $(1.5,\ 2.5)$ から広がるだけの流れ $\vec v_1$ と、口 $(4.5,\ 1.5)$ へ集まるだけの流れ $\vec v_2$ を、場所ごとにたしたものです（Python の式の \texttt{ax/a2 - bx/b2} の2つの項）。偏微分は、たし算の1つずつにできるので（和の微分 5-4 と同じ）、
\kviiijline{$\mathrm{div}(\vec v_1+\vec v_2)=\mathrm{div}\,\vec v_1+\mathrm{div}\,\vec v_2$}
$\vec v_1$ の発散は口 $(1.5,\ 2.5)$ 以外で0、$\vec v_2$ の発散は口 $(4.5,\ 1.5)$ 以外で0（⑪）。だから、池の流れの発散は、2つの口以外のどこでも0です。発散の地図で、0でないマスが2つだけだったのは、このためです。数をかけた流れも同じで、$\mathrm{div}(c\vec v)=c\,\mathrm{div}\,\vec v$（$c$ は数）。「たしてから発散を求めても、発散を求めてからたしても同じ」という性質は、7-2 の行列の「線形」と同じ考え方です。
\end{motto}

\begin{taisetsu}
\def\tsl#1{\makebox[9.6\zw][l]{\sffamily\bfseries\small\color{\bkTT}#1}}
\tsl{発散}小さな箱から出ていく量（出－入）÷箱の大きさ。答えは数\\[3pt]
\tsl{平面}$\mathrm{div}\,\vec v=\dfrac{\partial v_x}{\partial x}+\dfrac{\partial v_y}{\partial y}$\\[3pt]
\tsl{空間}$\mathrm{div}\,\vec v=\dfrac{\partial v_x}{\partial x}+\dfrac{\partial v_y}{\partial y}+\dfrac{\partial v_z}{\partial z}=\nabla\cdot\vec v$\\[3pt]
\tsl{マスの近い式}（右－左）÷$h$＋（上－下）÷$h$（右向き・上向きを正）\\[3pt]
\tsl{定義}$\displaystyle\lim_{\Delta V\to0}\frac{1}{\Delta V}\oiint_S\vec v\cdot\vec n\,\dd S$（8-8 の流束）\\[3pt]
\tsl{連続の式}$\dfrac{\partial\rho}{\partial t}+\mathrm{div}(\rho\vec v)=0$。縮まなければ $\mathrm{div}\,\vec v=0$
\end{taisetsu}



\bksection{taisetsu}{まとめ：発散を求める3つの方法}

\begin{bkcenter}
\begin{bktable}{colspec={Q[l,m,wd=18mm] Q[l,m,wd=62mm] Q[l,m,wd=22mm] Q[l,m,wd=16mm]}, rowsep=1.5pt, row{1}={font=\sffamily\small}}
段 & 方法 & 使う道具 & 例題 \\
目でわかる & マスを置いて、出る矢印と入る矢印をくらべる & 図・実験 & 1 \\
手でわかる & 4つの辺を通る量で「右－左＋上－下」を計算し、面積でわる & 表計算・Python & 2・3 \\
式でわかる & 偏微分の和 $\dfrac{\partial v_x}{\partial x}+\dfrac{\partial v_y}{\partial y}\Bigl(+\dfrac{\partial v_z}{\partial z}\Bigr)$ & 偏微分（8-2） & 4・5・7 \\
定義 & 小さな閉じた面の流束 $\oiint_S\vec v\cdot\vec n\,\dd S$ を体積でわる。実際には面ごとに計算（表8-8の \kviiijcase{G}） & 流束（8-8） & 6 \\
\end{bktable}
\bktabcaption{どの方法でも、同じ「発散」を求めている。マスや箱を小さくするほど、上の段の答えは下の段の答えに近づく。}
\end{bkcenter}

\begin{kotoba}
\small
\bkword{発散（divergence）}{小さな箱から出ていく量を、箱の大きさでわったもの。$\mathrm{div}\,\vec v$、$\nabla\cdot\vec v$ と書く。}
\bkword{わき出し・吸いこみ}{発散が正の所・負の所。物が生まれる所・消える所。}
\bkword{点のわき出し}{1点だけに集まったわき出し。その点では式で発散を計算できない（8-11）。}
\bkword{連続の式}{物が勝手に生まれも消えもしないことを、発散で書いた式。}
\bkword{非圧縮}{水のように縮まないこと。非圧縮の流れでは $\mathrm{div}\,\vec v=0$。}
\bkword{スカラー}{向きをもたず、大きさだけをもつ量。発散はスカラー。}
\end{kotoba}

\begin{senmonsho}[大学の教科書ふうの書き方]
\begin{bkhonbun}
\textbf{定義 2.5}\quad ベクトル場 $\symbf{A}(\symbf{r})=(A_x,\,A_y,\,A_z)$ に対して
\[ \operatorname{div}\symbf{A}\equiv\nabla\cdot\symbf{A}=\frac{\partial A_x}{\partial x}+\frac{\partial A_y}{\partial y}+\frac{\partial A_z}{\partial z} \tag{2.15} \]
を $\symbf{A}$ の発散という。
\end{bkhonbun}
\bkyaku{$\symbf{A}$ という矢印の場について、各成分を自分の向きの文字で偏微分して、たしたものを「発散」とよぶ。太字の $\symbf{A}$ は、この本の $\vec A$。$\equiv$ は「こう決める」という記号。（2.15）は、あとで引用するための式の番号。}
\begin{bkhonbun}
真空中のガウスの法則は、微分形で
\[ \nabla\cdot\symbf{E}=\frac{\rho}{\varepsilon_0} \tag{3.8} \]
と表される。
\end{bkhonbun}
\bkyaku{電場の発散（わき出し）は、その場所の電荷の密度 $\rho$ に比例する。「微分形」は、場所ごとの式（この章の形）という意味。袋全体で書いた「積分形」は 8-11 で。}
\begin{bkhonbun}
質量保存則より、連続の式
\[ \frac{\partial\rho}{\partial t}+\nabla\cdot(\rho\symbf{v})=0 \tag{4.2} \]
が成り立つ。特に非圧縮性流体では $\nabla\cdot\symbf{v}=0$ である。
\end{bkhonbun}
\bkyaku{「質量保存則」は、質量は生まれも消えもしないというきまり。⑫と同じ式。「非圧縮性流体」は、水のように縮まない流れる物。}
\end{senmonsho}





\bkfukidashi{delta}{専門書の式も、「小さな箱から出ていく量」と読めば、こわくないよ。}

%% ---------------- 例題 ----------------
\bksection{reidai}{例題でたしかめよう}

\begin{reidai}[図を読む]
3つの流れ（A）〜（C）で、点線のマスの発散は、正・負・0のどれでしょう。
\begin{bkfigure}
\begin{tabular}{@{}c@{\hspace{8mm}}c@{\hspace{8mm}}c@{}}
\kviiijpanel{-0.3*\x}{-0.3*\y}{\draw[kviiij box] (-0.5,-0.5) rectangle (0.5,0.5);} &
\kviiijpanel{0.5}{0.5}{\draw[kviiij box] (-0.5,-0.5) rectangle (0.5,0.5);} &
\kviiijpanel{0.16*(\x+3)}{0}{\draw[kviiij box] (-0.5,-0.5) rectangle (0.5,0.5);} \\
{\sffamily\bfseries\small（A）} & {\sffamily\bfseries\small（B）} & {\sffamily\bfseries\small（C）} \\
\end{tabular}
\end{bkfigure}
\kangaekata マスの辺を通って「出る量」と「入る量」をくらべる（①②）。矢印の向きではなく、出入りの差を見る。\par
\shiki\kviiijs\hangindent=2.6\zw\hangafter=1 （A）矢印がマスへ集まる。どの辺からも入る。入る量のほうが多い\par
\kviiija （B）ななめにそろった流れ。左と下の辺から入った分だけ、右と上の辺から出る\par
\kviiija （C）右へ行くほど矢印が長い。右の辺から出る量が、左の辺から入る量より多い\par
\kotae\kviiijs\hangindent=2.6\zw\hangafter=1 （A）負\quad （B）0\quad （C）正
\end{reidai}

\begin{jibun}
3つの流れ（D）〜（F）で、点線のマスの発散は、正・負・0のどれでしょう。
\begin{bkfigure}
\begin{tabular}{@{}c@{\hspace{8mm}}c@{\hspace{8mm}}c@{}}
\kviiijpanel{0.3*\x}{0.3*\y}{\draw[kviiij box] (-0.5,-0.5) rectangle (0.5,0.5);} &
\kviiijpanel{0.3*\y}{-0.3*\x}{\draw[kviiij box] (-0.5,-0.5) rectangle (0.5,0.5);} &
\kviiijpanel{-0.16*(\x+3)}{0}{\draw[kviiij box] (-0.5,-0.5) rectangle (0.5,0.5);} \\
{\sffamily\bfseries\small（D）} & {\sffamily\bfseries\small（E）} & {\sffamily\bfseries\small（F）} \\
\end{tabular}
\end{bkfigure}
\bkjans{（D）正（どの辺からも出る）。（E）0（うず：入った分だけ出る）。（F）負（矢印はみな左向きで、右へ行くほど長い。右の辺から入る量のほうが、左の辺から出る量より多い）。}
\end{jibun}

\begin{reidai}[マスの4つの辺から計算する]
1辺 $h=\qty{0.5}{m}$ のマスで、辺を通る水の速さ（右向き・上向きを正）が、次のようにわかっています。このマスの発散を求め、わき出しか吸いこみかを答えましょう。
\begin{bkfigure}
\begin{minipage}[c]{0.36\linewidth}\centering
\begin{tikzpicture}[x=1mm, y=1mm]
  \fill[bkOrange!8] (0,0) rectangle (22,22);
  \draw[kviiij box] (0,0) rectangle (22,22);
  \draw[bk arrow, bkTheme, line width=1.3pt] (23.5,11) -- (33.5,11);
  \node[font=\sffamily\footnotesize, text=bkInk, anchor=south] at (29,12.5) {1.0};
  \draw[bk arrow, bkTheme, line width=1.3pt] (-1.5,11) -- (-9.5,11);
  \node[font=\sffamily\footnotesize, text=bkInk, anchor=south] at (-5.5,12.5) {0.5};
  \draw[bk arrow, bkTheme, line width=1.3pt] (11,23.5) -- (11,29.5);
  \node[font=\sffamily\footnotesize, text=bkInk, anchor=west] at (12.5,27) {0.5};
  \draw[bk arrow, bkTheme, line width=1.3pt] (11,-12) -- (11,-1.5);
  \node[font=\sffamily\footnotesize, text=bkInk, anchor=west] at (12.5,-8) {1.0};
\end{tikzpicture}
\end{minipage}\hfill
\begin{minipage}[c]{0.6\linewidth}
\begin{tblr}{colspec={Q[l,m] Q[r,m] Q[l,m]}, rowsep=1pt, colsep=3pt}
右の辺 & $v_x=\qty{1.0}{m/s}$ & \\
左の辺 & $v_x=\qty{-0.5}{m/s}$ & （左向き） \\
上の辺 & $v_y=\qty{0.5}{m/s}$ & \\
下の辺 & $v_y=\qty{1.0}{m/s}$ & （上向き） \\
\end{tblr}
\end{minipage}
\end{bkfigure}
\kangaekata ④の式「（右－左）÷$h$＋（上－下）÷$h$」に入れる。左向きの速さはマイナスのまま入れる。\par
\shiki\kviiijs\hangindent=2.6\zw\hangafter=1 $\dfrac{1.0-(-0.5)}{0.5}+\dfrac{0.5-1.0}{0.5}$\par
\kviiija $=\dfrac{1.5}{0.5}+\dfrac{-0.5}{0.5}=3+(-1)=2$\par
\kotae\kviiijs\hangindent=2.6\zw\hangafter=1 $2$\,\kviiijps（正なので、わき出し）
\end{reidai}

\begin{reidai}[表計算で発散の地図を作る]
1辺\qty{1}{m}のマス4つ（2×2）で、辺を通る量（\unit{m^2/s}）が表のようにわかっています。4つのマスの発散と、いちばん強いわき出しのマスを答えましょう。また、4つの発散の合計と、外まわりから出る量をくらべましょう。
\begin{bkfigure}
\begin{minipage}[t]{0.5\linewidth}\centering
{\sffamily\small たての辺（右向きが正）}\par\vspace{2pt}
\begin{bktable}{colspec={Q[l,m,wd=12mm] Q[c,m,wd=11mm] Q[c,m,wd=11mm] Q[c,m,wd=11mm]}, rowsep=1pt, row{1}={font=\sffamily\footnotesize}}
段 & $x=0$ & $x=1$ & $x=2$ \\
上の段 & 0.5 & 1.0 & 0.5 \\
下の段 & 0.2 & 0.6 & 0.6 \\
\end{bktable}
\end{minipage}\hfill
\begin{minipage}[t]{0.48\linewidth}\centering
{\sffamily\small 横の辺（上向きが正）}\par\vspace{2pt}
\begin{bktable}{colspec={Q[l,m,wd=12mm] Q[c,m,wd=12mm] Q[c,m,wd=12mm]}, rowsep=1pt, row{1}={font=\sffamily\footnotesize}}
線 & 左の列 & 右の列 \\
$y=2$ & 0.3 & 0.4 \\
$y=1$ & 0.1 & $-0.2$ \\
$y=0$ & 0.1 & 0.0 \\
\end{bktable}
\end{minipage}
\end{bkfigure}
\kangaekata ⑤の表計算と同じ。マスごとに「右－左＋上－下」。面積は\qty{1}{m^2}なので、そのまま発散。上の段のマスの上の辺は $y=2$、下の辺は $y=1$ の線。\par
\shiki\kviiijs\hangindent=2.6\zw\hangafter=1 左上：$(1.0-0.5)+(0.3-0.1)=0.7$\par
\kviiija 右上：$(0.5-1.0)+(0.4-(-0.2))=0.1$\par
\kviiija 左下：$(0.6-0.2)+(0.1-0.1)=0.4$\par
\kviiija 右下：$(0.6-0.6)+(-0.2-0.0)=-0.2$\par
\kviiija 合計：$0.7+0.1+0.4+(-0.2)=1.0$\par
\kviiija 外まわり：右 $0.5+0.6$、左 $0.5+0.2$、上 $0.3+0.4$、下 $0.1+0.0$\par
\kviiijb 「右－左＋上－下」で $1.1-0.7+0.7-0.1=1.0$\par
\kotae\kviiijs\hangindent=2.6\zw\hangafter=1 左上 0.7、右上 0.1、左下 0.4、右下 $-0.2$（単位 \kviiijps）。いちばん強いわき出しは左上。合計も外まわりも 1.0 で等しい
\end{reidai}

\begingroup\tcbset{equal height group=kviiijEHI}%
\noindent\begin{minipage}[t]{0.485\linewidth}\vspace{0pt}
\begin{jibun}
\raggedright 1辺 $h=\qty{2}{m}$ のマスで、辺を通る水の速さ（右向き・上向きを正）は、右の辺 $v_x=\qty{-1}{m/s}$（左向き）、左の辺 $v_x=\qty{1}{m/s}$（右向き）、上の辺と下の辺はどちらも $v_y=\qty{1}{m/s}$（上向き）です。このマスの発散を求め、わき出しか吸いこみかを答えましょう。
\bkjans{$\dfrac{-1-1}{2}+\dfrac{1-1}{2}=-1+0=-1$。$-1$\,\kviiijps（負なので、吸いこみ）。}
\end{jibun}
\end{minipage}\hfill
\begin{minipage}[t]{0.485\linewidth}\vspace{0pt}
\begin{jibun}
\raggedright 例題3と同じ並びの4つのマスで、たての辺（右向きが正）は上の段が 0.4、0.2、0.3、下の段が 0.0、0.5、0.1、横の辺（上向きが正）は $y=2$ が 0.2、0.1、$y=1$ が 0.3、$-0.1$、$y=0$ が $-0.1$、0.2 です（\unit{m^2/s}）。4つのマスの発散と、いちばん強いわき出しのマスを答え、発散の合計と外まわりから出る量をくらべましょう。
\bkjans{左上 $(0.2-0.4)+(0.2-0.3)=-0.3$、右上 $(0.3-0.2)+(0.1-(-0.1))=0.3$、左下 $(0.5-0.0)+(0.3-(-0.1))=0.9$、右下 $(0.1-0.5)+(-0.1-0.2)=-0.7$（単位 \kviiijps）。いちばん強いわき出しは左下。合計 $0.2$。外まわりは、右 $0.3+0.1$、左 $0.4+0.0$、上 $0.2+0.1$、下 $-0.1+0.2$ から $0.4-0.4+0.3-0.1=0.2$ で等しい。}
\end{jibun}
\end{minipage}\par
\endgroup

\begin{reidai}[式で表された平面の流れ]
（1）平面の流れ $\vec v=(3x,\ -y)$ の発散を求めましょう。\par
（2）流れ $\vec v=(ax,\ -y)$ が、どこでも発散0になるのは、$a$ がいくつのときでしょう。\par
\kangaekata ⑥の式。$v_x$ は $x$ で、$v_y$ は $y$ で偏微分する。\par
\shiki\kviiijs\hangindent=2.6\zw\hangafter=1 （1）$\dfrac{\partial v_x}{\partial x}=\dfrac{\partial(3x)}{\partial x}=3$、$\dfrac{\partial v_y}{\partial y}=\dfrac{\partial(-y)}{\partial y}=-1$\par
\kviiijb $\mathrm{div}\,\vec v=3+(-1)=2$\par
\kviiija （2）$\mathrm{div}\,\vec v=\dfrac{\partial(ax)}{\partial x}+\dfrac{\partial(-y)}{\partial y}=a-1$\par
\kviiijb $a-1=0$ より $a=1$\par
\kotae\kviiijs\hangindent=2.6\zw\hangafter=1 （1）$\mathrm{div}\,\vec v=2$（どこでも正なので、どこでもわき出し）\quad （2）$a=1$
\end{reidai}

\begin{reidai}[場所によって変わる発散（空間）]
空間の流れ $\vec v=(x^2,\ xy,\ -z)$ について答えましょう。\par
（1）発散を求めましょう。\par
（2）点 $(1,\ 0,\ 0)$ と点 $(0,\ 2,\ 5)$ は、わき出しと吸いこみのどちらでしょう。\par
（3）発散が0になる点は、どんな所に集まっていますか。\par
\kangaekata ⑧の式。微分する文字以外は、数と同じに扱う（$\frac{\partial(xy)}{\partial y}$ では $x$ を止めるので $x$）。\par
\shiki\kviiijs\hangindent=2.6\zw\hangafter=1 （1）$\dfrac{\partial(x^2)}{\partial x}=2x$、$\dfrac{\partial(xy)}{\partial y}=x$、$\dfrac{\partial(-z)}{\partial z}=-1$\par
\kviiijb $\mathrm{div}\,\vec v=2x+x-1=3x-1$\par
\kviiija （2）点 $(1,\ 0,\ 0)$：$3\times1-1=2>0$\par
\kviiijb 点 $(0,\ 2,\ 5)$：$3\times0-1=-1<0$\par
\kviiija （3）$3x-1=0$ より $x=\frac13$（$y$・$z$ は何でもよい）\par
\kotae\kviiijs\hangindent=2.6\zw\hangafter=1 （1）$3x-1$\quad （2）$(1,\ 0,\ 0)$ はわき出し、$(0,\ 2,\ 5)$ は吸いこみ\par
\kviiija （3）$x=\frac13$ の平面の上（その右側はわき出し、左側は吸いこみ）
\end{reidai}

\begingroup\tcbset{equal height group=kviiijEHJ}%
\noindent\begin{minipage}[t]{0.485\linewidth}\vspace{0pt}
\begin{jibun}
\raggedright （1）平面の流れ $\vec v=(x+2y,\ 4x-3y)$ の発散を求めましょう。\par
（2）平面の流れ $\vec v=(2x,\ by)$ が、どこでも発散0になるのは、$b$ がいくつのときでしょう。
\bkjans{（1）$\dfrac{\partial(x+2y)}{\partial x}+\dfrac{\partial(4x-3y)}{\partial y}=1+(-3)=-2$（どこでも吸いこみ）。（2）$2+b=0$ より $b=-2$。}
\end{jibun}
\end{minipage}\hfill
\begin{minipage}[t]{0.485\linewidth}\vspace{0pt}
\begin{jibun}
\raggedright 空間の流れ $\vec v=(xy,\ y^2,\ z)$ について答えましょう。\par
（1）発散を求めましょう。\par
（2）点 $(2,\ 1,\ 0)$ と点 $(0,\ -1,\ 3)$ は、わき出しと吸いこみのどちらですか。\par
（3）発散が0になる点は、どんな所に集まっていますか。
\bkjans{（1）$y+2y+1=3y+1$。（2）$(2,\ 1,\ 0)$ は $3+1=4>0$ でわき出し、$(0,\ -1,\ 3)$ は $-3+1=-2<0$ で吸いこみ。（3）$y=-\frac13$ の平面の上。}
\end{jibun}
\end{minipage}\par
\endgroup

\begin{reidai}[小さな箱の面ごとの流束から求める]
空間の流れ $\vec v=(x^2,\ xy,\ -z)$ の中で、点 $\mathrm P(1,\ 0,\ 0)$ を中心に、1辺 $h$ の立方体の箱を、面が軸に直角になるように置きます。6つの面から外へ出る量を面ごとに求めてたし、体積 $h^3$ でわって、点Pの発散を求めましょう。例題5の答えとくらべましょう。\par
\kangaekata \kviiijbaai{閉じた面が立方体の6つの平らな面}{8-8}{G}\par
\kviiijdan{1}\enspace $\displaystyle\oiint_S\vec v\cdot\vec n\,\dd S$：箱の面全体で「外向きの成分×タイル」をたす。\par
\kviiijdan{3}\enspace 面ごとに分ける。成分が面の上で1次式で変わるときは「面のまん中の値×面積 $h^2$」でよい（平均がまん中の値になるから）。\par
\kviiijdan{2}\enspace 6つの面それぞれで、$\vec n$ と外向きの成分を表にする。\par
\begin{bkcenter}
\begin{bktable}{colspec={Q[l,m,wd=31mm] Q[c,m,wd=19mm] Q[c,m,wd=21mm] Q[c,m,wd=29mm]}, rowsep=0.5pt, colsep=4pt, row{1}={font=\sffamily\footnotesize}}
面 & $\vec n$ & 外向きの成分 & 面のまん中の値 \\
右（$x=1+\frac h2$） & $(1,\ 0,\ 0)$ & $v_x=x^2$ & $\bigl(1+\frac h2\bigr)^2$ \\
左（$x=1-\frac h2$） & $(-1,\ 0,\ 0)$ & $-v_x$ & $-\bigl(1-\frac h2\bigr)^2$ \\
奥（$y=\frac h2$） & $(0,\ 1,\ 0)$ & $v_y=xy$ & $1\times\frac h2=\frac h2$ \\
手前（$y=-\frac h2$） & $(0,\ -1,\ 0)$ & $-v_y$ & $-1\times\bigl(-\frac h2\bigr)=\frac h2$ \\
上（$z=\frac h2$） & $(0,\ 0,\ 1)$ & $v_z=-z$ & $-\frac h2$ \\
下（$z=-\frac h2$） & $(0,\ 0,\ -1)$ & $-v_z=z$ & $-\frac h2$ \\
\end{bktable}
\end{bkcenter}
\shiki\kviiijs\hangindent=2.6\zw\hangafter=1 左右：$\bigl\{\bigl(1+\frac h2\bigr)^2-\bigl(1-\frac h2\bigr)^2\bigr\}\times h^2=2h\times h^2=2h^3$\par
\kviiija 奥と手前：$\bigl(\frac h2+\frac h2\bigr)\times h^2=h^3$\par
\kviiija 上下：$\bigl(-\frac h2-\frac h2\bigr)\times h^2=-h^3$\par
\kviiija 合計 $2h^3+h^3-h^3=2h^3$。体積でわって $\dfrac{2h^3}{h^3}=2$\par
\kotae\kviiijs\hangindent=2.6\zw\hangafter=1 2。例題5の $3x-1$ に $x=1$ を入れた2と同じ
\end{reidai}
\noindent 例題6の左右の $\bigl(1+\frac h2\bigr)^2-\bigl(1-\frac h2\bigr)^2$ は、$(a+b)^2-(a-b)^2=4ab$ を使うと $4\times1\times\frac h2=2h$ です。この場では、$h$ がどんな大きさでも答えが2になりました。場が2次式までなら、箱の大きさによらず、ちょうどの値が出ます。

\begin{jibun}
空間の流れ $\vec v=(xy,\ y^2,\ z)$ の中で、点 $\mathrm P(2,\ 1,\ 0)$ を中心に、1辺 $h$ の立方体の箱を置きます。6つの面から外へ出る量を面ごとに求めてたし、体積 $h^3$ でわって、点Pの発散を求めましょう。
\bkjans{左右：$v_x=xy$ のまん中の値は $\bigl(2+\frac h2\bigr)\times1$ と $-\bigl(2-\frac h2\bigr)\times1$。たして $h$、面積をかけて $h^3$。奥と手前：$\bigl(1+\frac h2\bigr)^2-\bigl(1-\frac h2\bigr)^2=2h$ で $2h^3$。上下：$\frac h2-\bigl(-\frac h2\bigr)=h$ で $h^3$。合計 $4h^3$、$h^3$ でわって4（$3y+1$ に $y=1$ を入れた値と同じ）。}
\end{jibun}
\begin{reidai}[点のわき出しの場の発散]
平面の流れ $\vec v=\dfrac{(x,\ y)}{x^2+y^2}$ の発散を、原点以外の点で求めましょう。また、点 $(1,\ 2)$ での $\dfrac{\partial v_x}{\partial x}$ と $\dfrac{\partial v_y}{\partial y}$ の値を求めましょう。\par
\kangaekata ⑪と同じ。商の微分（5-6）で、止める文字に気をつける。\par
\shiki\kviiijs\hangindent=2.6\zw\hangafter=1 $\dfrac{\partial}{\partial x}\Bigl(\dfrac{x}{x^2+y^2}\Bigr)=\dfrac{(x^2+y^2)-x\times2x}{(x^2+y^2)^2}=\dfrac{y^2-x^2}{(x^2+y^2)^2}$\par
\kviiija $\dfrac{\partial}{\partial y}\Bigl(\dfrac{y}{x^2+y^2}\Bigr)=\dfrac{(x^2+y^2)-y\times2y}{(x^2+y^2)^2}=\dfrac{x^2-y^2}{(x^2+y^2)^2}$\par
\kviiija たして $\dfrac{(y^2-x^2)+(x^2-y^2)}{(x^2+y^2)^2}=0$\par
\kviiija 点 $(1,\ 2)$：$\dfrac{4-1}{5^2}=\dfrac{3}{25}=0.12$、$\dfrac{1-4}{5^2}=-0.12$\par
\kotae\kviiijs\hangindent=2.6\zw\hangafter=1 原点以外で0。点 $(1,\ 2)$ では $0.12$ と $-0.12$（たすと0）
\end{reidai}

\begin{reidai}[ホースの先：縮まない流れ]
庭のホースの中を、水が速さ\qty{1.5}{m/s}で流れています。ホースの断面積は\qty{2.0}{cm^2}です。先を指でつまんで、出口の断面積を\qty{0.5}{cm^2}にしました。\par
（1）出口での水の速さを求めましょう。\par
（2）1分間に出る水は何Lでしょう。\par
\kangaekata \kviiijbaai{管の断面は平らで、流れはそろっている}{8-8}{A}\par
水は縮まないので $\mathrm{div}\,\vec v=0$。袋の中で生まれも消えもしないので、入る量＝出る量（⑫）。断面を通る量＝速さ×断面積。単位をそろえる：$\qty{1}{cm^2}=\qty{0.0001}{m^2}$。\par
\shiki\kviiijs\hangindent=2.6\zw\hangafter=1 （1）$v_1A_1=v_2A_2$ より $1.5\times2.0=v_2\times0.5$\par
\kviiijb $v_2=\dfrac{1.5\times2.0}{0.5}=6.0$\par
\kviiija （2）1秒間：$\qty{1.5}{m/s}\times\qty{0.0002}{m^2}=\qty{0.0003}{m^3/s}$\par
\kviiijb 1分間：$\qty{0.0003}{m^3/s}\times\qty{60}{s/min}=\qty{0.018}{m^3/min}$\par
\kviiijb $\qty{0.018}{m^3/min}\times\qty{1000}{L/m^3}=\qty{18}{L/min}$\par
\kotae\kviiijs\hangindent=2.6\zw\hangafter=1 （1）\qty{6.0}{m/s}（断面積が $\frac14$ なので4倍）\quad （2）\qty{18}{L}
\end{reidai}

\noindent （1）の式では、両辺の断面積を cm² のまま使いました。両辺に同じ単位が出てくるので、m² に直さなくても比は同じだからです。（2）は量そのものを求めるので、m² に直しました。

\begingroup\tcbset{equal height group=kviiijEHK}%
\noindent\begin{minipage}[t]{0.485\linewidth}\vspace{0pt}
\begin{jibun}
\raggedright 平面の流れ $\vec v=\dfrac{(-2x,\ -2y)}{x^2+y^2}$（原点に吸いこみ口がある流れ）の発散を、原点以外の点で求めましょう。また、点 $(2,\ 1)$ での $\dfrac{\partial v_x}{\partial x}$ と $\dfrac{\partial v_y}{\partial y}$ の値を求めましょう。
\bkjans{$\dfrac{\partial v_x}{\partial x}=-2\times\dfrac{y^2-x^2}{(x^2+y^2)^2}$、$\dfrac{\partial v_y}{\partial y}=-2\times\dfrac{x^2-y^2}{(x^2+y^2)^2}$ で、たすと0（原点以外）。点 $(2,\ 1)$ では $-2\times\dfrac{1-4}{25}=0.24$ と $-2\times\dfrac{4-1}{25}=-0.24$。}
\end{jibun}
\end{minipage}\hfill
\begin{minipage}[t]{0.485\linewidth}\vspace{0pt}
\begin{jibun}
\raggedright ホースの中を、水が速さ\qty{0.8}{m/s}で流れています。ホースの断面積は\qty{2.5}{cm^2}、ノズルの出口の断面積は\qty{0.4}{cm^2}です。（1）出口での水の速さ、（2）1分間に出る水の量（L）を求めましょう。
\bkjans{（1）$0.8\times2.5=v_2\times0.4$ より $v_2=\qty{5.0}{m/s}$。（2）$\qty{0.8}{m/s}\times\qty{0.00025}{m^2}=\qty{0.0002}{m^3/s}$、$\times\qty{60}{s/min}=\qty{0.012}{m^3/min}$、$\times\qty{1000}{L/m^3}=\qty{12}{L}$。}
\end{jibun}
\end{minipage}\par
\endgroup


\begin{reidai}[熱の流れの発散：生まれる熱]
電気を流して熱を出している金属のブロックの中で、熱の流れの矢印が $\vec q=(2000x,\ 2000y,\ 2000z)$（W/m²）でした（$x$・$y$・$z$ はブロックのまん中からの場所〔m〕）。\par
（1）$\mathrm{div}\,\vec q$ を求め、単位をつけて答えましょう。\par
（2）ブロックの体積は\qty{0.5}{L}です。ブロック全体で、1秒間に生まれる熱は何Wでしょう。\par
\kangaekata ⑬。$\mathrm{div}\,\vec q$ は「\qty{1}{m^3}あたり、1秒間に出ていく熱」。出ていく熱は、その場で生まれた熱。偏微分は「÷長さ」なので、単位は $\unit{W/m^2}\div\unit{m}=\unit{W/m^3}$。\par
\shiki\kviiijs\hangindent=2.6\zw\hangafter=1 （1）$\dfrac{\partial(2000x)}{\partial x}+\dfrac{\partial(2000y)}{\partial y}+\dfrac{\partial(2000z)}{\partial z}=2000\times3=6000$\par
\kviiija （2）$\qty{0.5}{L}=\qty{0.0005}{m^3}$（$\qty{1}{L}=\qty{0.001}{m^3}$）\par
\kviiijb どこでも同じ $\qty{6000}{W/m^3}$ なので、$\qty{6000}{W/m^3}\times\qty{0.0005}{m^3}=\qty{3}{W}$\par
\kotae\kviiijs\hangindent=2.6\zw\hangafter=1 （1）\qty{6000}{W/m^3}（どこでも熱が生まれている）\quad （2）\qty{3}{W}
\end{reidai}

\begin{jibun}
電気を流して熱を出している板の中で、熱の流れの矢印 $\vec q$〔W/m²〕の成分は、$q_x=1500x$、$q_y=1500y$、$q_z=0$ でした。（1）$\mathrm{div}\,\vec q$ を単位をつけて求め、（2）板の体積が\qty{2}{L}のとき、板全体で1秒間に生まれる熱は何Wか答えましょう。
\bkjans{（1）$1500+1500+0=\qty{3000}{W/m^3}$。（2）$\qty{2}{L}=\qty{0.002}{m^3}$、$\qty{3000}{W/m^3}\times\qty{0.002}{m^3}=\qty{6}{W}$。}
\end{jibun}

%% ---------------- まちがえやすいところ ----------------
\bksection{tsumazuki}{まちがえやすいところをチェック}

\begin{tsumazuki}{① 広がる矢印なら、どこでもわき出し？}
\bkng{1つの口から水がわき出す池では、矢印はどこでも外向きなので、どこでも発散は正}
\bkok{口からはなれた所では、発散は0（図\ref{bk:8-9:fig-ten}の点線の箱）}
外へ行くほど矢印は短くなりますが、そのぶん外側の辺が長くなり、出る量と入る量は等しくなります。わき出しは口の1点だけです。発散は、矢印の向きではなく、出入りの差で決まります。
\end{tsumazuki}

\begin{tsumazuki}{② 成分と文字の組み合わせをまちがえる}
\bkng{$\vec v=(y,\ x)$ の発散を $\dfrac{\partial v_x}{\partial y}+\dfrac{\partial v_y}{\partial x}=1+1=2$}
\bkok{$\dfrac{\partial v_x}{\partial x}+\dfrac{\partial v_y}{\partial y}=0+0=0$}
左右の出入りは $v_x$ の左右の差、上下の出入りは $v_y$ の上下の差で決まります。だから $v_x$ は $x$ で、$v_y$ は $y$ で微分します。$\nabla\cdot\vec v$ の「1つめどうし」と覚えてもよいです。
\end{tsumazuki}

\begingroup\tcbset{equal height group=kviiijEHT}%
\noindent\begin{minipage}[t]{0.485\linewidth}\vspace{0pt}
\begin{tsumazuki}{③ 発散を矢印で答える}
\raggedright\bkng{$\vec v=(3x,\ -y)$ の発散は $(3,\ -1)$}
\bkok{$\mathrm{div}\,\vec v=3+(-1)=2$}
発散は、場所ごとに1つの数（スカラー）です。成分ごとの偏微分を、最後に必ずたします。矢印が出てくるのは grad（8-4）と rot（8-10）です。
\end{tsumazuki}
\end{minipage}\hfill
\begin{minipage}[t]{0.485\linewidth}\vspace{0pt}
\begin{tsumazuki}{④ 左と下の辺も、たしてしまう}
\raggedright\bkng{右 3、左 1、上 2、下 2（どれも右向き・上向き）で、$3+1+2+2=8$}
\bkok{$(3-1)+(2-2)=2$}
左の辺で右向きに流れる水は、マスに「入る」水です。右向き・上向きを正にしたら、左と下の辺は\textbf{ひく}（表\ref{bk:8-9:tab-fugou}）。
\end{tsumazuki}
\end{minipage}\par
\endgroup

\begingroup\tcbset{equal height group=kviiijEHA}%
\noindent\begin{minipage}[t]{0.485\linewidth}\vspace{0pt}
\begin{tsumazuki}{⑤ 口の外が0なら、わき出しはない？}
\raggedright\bkng{点のわき出しの場は、式で発散が0なので、水はどこからもわいていない}
\bkok{原点（口）の1点に、わき出し $Q$ が集まっている}
原点では式が使えません。口を囲めば、外へ出る量は $Q$ です（8-11）。
\end{tsumazuki}
\end{minipage}\hfill
\begin{minipage}[t]{0.485\linewidth}\vspace{0pt}
\begin{tsumazuki}{⑥ $\oiint$ の◯を球だと思う}
\raggedright\bkng{定義の $\oiint_S$ は、球の面でしか使えない}
\bkok{◯は「閉じている」という印。立方体でも球でもよい}
どの面かは下の $S$ で決めます。小さくすれば、どの形でも同じ発散になります。
\end{tsumazuki}
\end{minipage}\par
\endgroup

\begin{column}{天気予報の「発散」と「収束」}
\noindent\begin{minipage}[c]{0.6\linewidth}
\raggedright
天気予報では、地面の近くの風が外へ広がる所を「発散」、集まる所を「収束」とよびます。高気圧の中心付近の地面近くでは、風が外へ吹き出し（発散）、上から空気が下りてきて、晴れやすくなります。低気圧では地面近くの風が集まり（収束）、行き場のない空気が上へのぼって、雲ができやすくなります。空気の連続の式（⑫）で、地面近くの発散・収束と、上下の空気の動きが結びついているのです。
\end{minipage}\hfill
\begin{minipage}[c]{0.37\linewidth}
\centering
\begin{tikzpicture}[x=1mm, y=1mm]
  % 高気圧
  \fill[bkOrange!15] (0,0) ellipse (13 and 4);
  \foreach \a in {0,60,...,300} {\draw[bk arrow, bkTheme, line width=0.9pt] (\a:3 and 1) -- (\a:12 and 3.7);}
  \draw[bk arrow, bkTheme, line width=0.9pt] (0,16) -- (0,6);
  \node[font=\sffamily\footnotesize, text=bkInk] at (0,-8) {高気圧（発散）};
  % 低気圧
  \begin{scope}[shift={(0,-40)}]
  \fill[bkBlue!15] (0,0) ellipse (13 and 4);
  \foreach \a in {0,60,...,300} {\draw[bk arrow, bkTheme, line width=0.9pt] (\a:12 and 3.7) -- (\a:3 and 1);}
  \draw[bk arrow, bkTheme, line width=0.9pt] (0,5) -- (0,15);
  \fill[bkGray!30] (-5,19) circle (3) (0,20.5) circle (3.6) (5,19) circle (3);
  \node[font=\sffamily\footnotesize, text=bkInk] at (0,-8) {低気圧（収束）};
  \end{scope}
\end{tikzpicture}
\end{minipage}
\end{column}

%% ---------------- 物理・工学ではこう使う ----------------
\begin{butsuri}
発散は、「何かが生まれたり、消えたりしている場所」を見つける道具です。
\begin{bkfigure}
\begin{tikzpicture}[x=1mm, y=1mm]
  % 電荷
  \begin{scope}
    \fill[white, rounded corners=2mm] (0,0) rectangle (38,32);
    \fill[bkRed] (11,19) circle (2.6); \node[text=white, font=\bfseries] at (11,19) {＋};
    \fill[bkBlue] (27,19) circle (2.6); \node[text=white, font=\bfseries] at (27,19) {－};
    \foreach \a in {0,45,...,315} {\draw[bk arrow, bkRed!80!black, line width=0.8pt] (11,19) ++(\a:3.6) -- ++(\a:5);}
    \foreach \a in {0,45,...,315} {\draw[bk arrow, bkBlue!80!black, line width=0.8pt] (27,19) ++(\a:9) -- ++(\a:-5);}
    \node[font=\sffamily\bfseries\footnotesize, text=bkInk] at (19,4) {電荷と電場};
  \end{scope}
  % 磁石
  \begin{scope}[shift={(42,0)}]
    \fill[white, rounded corners=2mm] (0,0) rectangle (38,32);
    \fill[bkRed] (11,17) rectangle (19,22); \fill[bkBlue] (19,17) rectangle (27,22);
    \node[text=white, font=\sffamily\bfseries\footnotesize] at (15,19.5) {N}; \node[text=white, font=\sffamily\bfseries\footnotesize] at (23,19.5) {S};
    \foreach \s in {1,1.8} {\draw[bkPurple!80!black, line width=0.7pt] (11,19.5) .. controls (11-6*\s,19.5+8*\s) and (27+6*\s,19.5+8*\s) .. (27,19.5);
      \draw[bkPurple!80!black, line width=0.7pt] (11,19.5) .. controls (11-6*\s,19.5-8*\s) and (27+6*\s,19.5-8*\s) .. (27,19.5);}
    \node[font=\sffamily\bfseries\footnotesize, text=bkInk] at (19,4) {磁石と磁場};
  \end{scope}
  % 水道管
  \begin{scope}[shift={(84,0)}]
    \fill[white, rounded corners=2mm] (0,0) rectangle (38,32);
    \fill[bkBlueBg] (3,12) -- (16,12) -- (25,15) -- (35,15) -- (35,23) -- (25,23) -- (16,26) -- (3,26) -- cycle;
    \draw[bkGray, line width=0.8pt] (3,12) -- (16,12) -- (25,15) -- (35,15) (3,26) -- (16,26) -- (25,23) -- (35,23);
    \draw[bk arrow, bkTheme, line width=1pt] (5,19) -- (11,19);
    \draw[bk arrow, bkTheme, line width=1pt] (25,19) -- (35,19);
    \node[font=\sffamily\bfseries\footnotesize, text=bkInk] at (19,4) {細い所で速くなる};
  \end{scope}
\end{tikzpicture}
\end{bkfigure}

\kviiijbox{電気と磁気}
電場の発散は電荷の密度に比例し、プラスの電荷はわき出し、マイナスの電荷は吸いこみです（$\mathrm{div}\,\vec E=\rho/\varepsilon_0$）。半導体の中の電場や電位の計算（トランジスタの設計）も、この式から始まります。磁場には、わき出しも吸いこみもありません（$\mathrm{div}\,\vec B=0$）。磁力線は、いつも輪になっています。

\kviiijbox{水の流れ}
水はほとんど縮まないので $\mathrm{div}\,\vec v=0$。管が細い所で水が速くなるのは、このためです。水道管・ポンプ・ノズルの設計や、コンピュータで流れを計算するとき（流体の数値計算）は、「どの小さな箱でも、入る量＝出る量」になるように計算を組み立てます。

\kviiijbox{熱}
熱の流れの発散は、そこから出ていく熱の割合です。電子機器の部品のように熱が生まれる所で正になります。基板の設計では、熱がどこで生まれ、どこへ流れ出るかを、この考えで計算します（8-12・8-13）。

\end{butsuri}

%% ---------------- 章末問題 ----------------
\bkendproblems
\bkproblemgroup{基本}{5}

\bkproblem 3つの流れ（ア）〜（ウ）で、点線のマスの発散は、正・負・0のどれでしょう\mbox{（\bkreidairef{1}）}。
\begin{bkfigure}
\begin{tabular}{@{}c@{\hspace{8mm}}c@{\hspace{8mm}}c@{}}
\kviiijpanel{0}{-0.6}{\draw[kviiij box] (-0.5,-0.5) rectangle (0.5,0.5);} &
\kviiijpanel{0}{0.14*(\y+3)}{\draw[kviiij box] (-0.5,-0.5) rectangle (0.5,0.5);} &
\kviiijpanel{-0.3*\x}{-0.3*\y}{\draw[kviiij box] (0.5,-0.5) rectangle (1.5,0.5);} \\
{\sffamily\bfseries\small（ア）} & {\sffamily\bfseries\small（イ）} & {\sffamily\bfseries\small（ウ）} \\
\end{tabular}
\end{bkfigure}
\bkpans{（ア）0（下向きにそろった流れ）。（イ）正（上へ行くほど速い：上の辺から出る量が多い）。（ウ）負（集まる流れは、どこに置いたマスでも入る量が多い）。}

\bkproblem 次のマスの発散を求め、わき出しか吸いこみかを答えましょう（右向き・上向きを正）\mbox{（\bkreidairef{2}）}。\par
（1）1辺 $h=\qty{1}{m}$ のマスで、右の辺 $v_x=\qty{4}{m/s}$、左の辺 $v_x=\qty{1}{m/s}$、上の辺 $v_y=\qty{0}{m/s}$、下の辺 $v_y=\qty{2}{m/s}$\par
（2）1辺 $h=\qty{0.2}{m}$ のマスで、右の辺 $v_x=\qty{0.3}{m/s}$、左の辺 $v_x=\qty{0.5}{m/s}$、上の辺 $v_y=\qty{-0.2}{m/s}$、下の辺 $v_y=\qty{0.1}{m/s}$
\bkpans{（1）$\dfrac{4-1}{1}+\dfrac{0-2}{1}=3-2=1$。$1$\,\kviiijps、わき出し。（2）$\dfrac{0.3-0.5}{0.2}+\dfrac{-0.2-0.1}{0.2}=-1+(-1.5)=-2.5$。$-2.5$\,\kviiijps、吸いこみ。}

\bkproblem 例題3と同じ並びの4つのマスで、たての辺を右向きに通る量は、上の段が $x=0$ から順に 0.1、0.6、0.2、下の段が 0.3、0.3、0.5 です。横の辺を上向きに通る量は、$y=2$ の線が左の列から 0.4、0.0、$y=1$ の線が 0.2、0.1、$y=0$ の線が 0.0、0.2 です（単位は \unit{m^2/s}）。4つのマスの発散を求めましょう。また、4つの発散の合計と、外まわりから出る量をくらべましょう\mbox{（\bkreidairef{3}）}。
\bkpans{左上 $(0.6-0.1)+(0.4-0.2)=0.7$、右上 $(0.2-0.6)+(0.0-0.1)=-0.5$、左下 $(0.3-0.3)+(0.2-0.0)=0.2$、右下 $(0.5-0.3)+(0.1-0.2)=0.1$（単位 \kviiijps）。合計 0.5。外まわりは、右 $0.2+0.5$、左 $0.1+0.3$、上 $0.4+0.0$、下 $0.0+0.2$ から $0.7-0.4+0.4-0.2=0.5$ で等しい。}

\bkproblem 次の平面の流れの発散を求めましょう\mbox{（\bkreidairef{4}）}。\par
（1）$\vec v=(5x,\ 2y)$\par
（2）$\vec v=(y^2,\ x^2)$\par
（3）$\vec v=(2x-y,\ x+y)$
\bkpans{（1）$5+2=7$。（2）$0+0=0$。（3）$2+1=3$。}

\bkproblem 平面の流れ $\vec v=(4x+y,\ cy-x)$ が、どこでも発散0（縮まない流れ）になるのは、$c$ がいくつのときでしょう\mbox{（\bkreidairef{4}）}。
\bkpans{$\mathrm{div}\,\vec v=4+c=0$ より $c=-4$。}

\bkproblem 空間の流れ $\vec v=(xz,\ 2y,\ -z^2)$ について答えましょう\mbox{（\bkreidairef{5}）}。\par
（1）発散を求めましょう。\par
（2）点 $(1,\ 0,\ 1)$ と点 $(0,\ 5,\ 3)$ は、わき出しと吸いこみのどちらでしょう。\par
（3）発散が0になる点は、どんな所に集まっていますか。
\bkpans{（1）$z+2-2z=2-z$。（2）$(1,\ 0,\ 1)$ は $2-1=1>0$ でわき出し、$(0,\ 5,\ 3)$ は $2-3=-1<0$ で吸いこみ。（3）$z=2$ の平面の上。}

\bkproblem ホースの中を、水が速さ\qty{1.2}{m/s}で流れています。ホースの断面積は\qty{3.0}{cm^2}、ノズルの出口の断面積は\qty{0.6}{cm^2}です。（1）出口での水の速さ、（2）1分間に出る水の量（L）を求めましょう\mbox{（\bkreidairef{8}）}。
\bkpans{（1）$1.2\times3.0=v_2\times0.6$ より $v_2=\qty{6.0}{m/s}$。（2）$\qty{1.2}{m/s}\times\qty{0.0003}{m^2}=\qty{0.00036}{m^3/s}$、$\times\qty{60}{s/min}=\qty{0.0216}{m^3/min}$、$\times\qty{1000}{L/m^3}=\qty{21.6}{L}$。}

\bkproblem 電気を流して熱を出している部品の中で、熱の流れの矢印 $\vec q$〔W/m²〕の成分は、$q_x=800x$、$q_y=800y$、$q_z=800z$ でした。（1）$\mathrm{div}\,\vec q$ を単位をつけて求め、（2）部品の体積が\qty{1.5}{L}のとき、部品全体で1秒間に生まれる熱は何Wか答えましょう\mbox{（\bkreidairef{9}）}。
\bkpans{（1）$800\times3=\qty{2400}{W/m^3}$。（2）$\qty{1.5}{L}=\qty{0.0015}{m^3}$、$\qty{2400}{W/m^3}\times\qty{0.0015}{m^3}=\qty{3.6}{W}$。}

\bkproblemgroup{チャレンジ}{5}

\bkproblem 平面の流れ $\vec v=(x^2,\ -2xy)$ について答えましょう\mbox{（\bkreidairef{4}）}。\par
（1）発散を求めましょう。\par
（2）この流れの中に、わき出しや吸いこみはあるでしょうか。
\bkpans{（1）$2x+(-2x)=0$。（2）どこにもない（どこでも発散0の、縮まない流れ）。}

\bkproblem 空間の流れ $\vec v=(xy,\ yz,\ zx)$ について答えましょう\mbox{（\bkreidairef{5}）}。\par
（1）発散を求めましょう。\par
（2）点 $(1,\ 2,\ 3)$ での発散の値を求めましょう。\par
（3）点 $(1,\ -2,\ 0)$ は、わき出しと吸いこみのどちらでしょう。
\bkpans{（1）$y+z+x$。（2）$2+3+1=6$。（3）$-2+0+1=-1<0$ で吸いこみ。}

\bkproblem 空間の流れ $\vec v=(x^2,\ xy,\ -z)$ の中で、点 $(2,\ 0,\ 0)$ を中心に1辺 $h$ の立方体の箱を置きます。6つの面から外へ出る量を面ごとに求めてたし、体積 $h^3$ でわって、この点の発散を求めましょう。例題5の式の値とくらべましょう\mbox{（\bkreidairef{6}）}。
\bkpans{左右：$\bigl(2+\frac h2\bigr)^2-\bigl(2-\frac h2\bigr)^2=4h$ で $4h^3$。奥と手前：$v_y=xy$ のまん中の値は $2\times\frac h2$ と $-2\times\bigl(-\frac h2\bigr)$、たして $2h$ で $2h^3$。上下：$-\frac h2-\frac h2=-h$ で $-h^3$。合計 $5h^3$、$h^3$ でわって5（$3x-1$ に $x=2$ を入れた値と同じ）。}

\bkproblem 平面の流れ $\vec v=\dfrac{(3x,\ 3y)}{x^2+y^2}$ の発散を、原点以外の点で求めましょう。また、点 $(1,\ 2)$ での $\dfrac{\partial v_x}{\partial x}$ と $\dfrac{\partial v_y}{\partial y}$ の値を求めましょう\mbox{（\bkreidairef{7}）}。
\bkpans{$\dfrac{\partial v_x}{\partial x}=3\times\dfrac{y^2-x^2}{(x^2+y^2)^2}$、$\dfrac{\partial v_y}{\partial y}=3\times\dfrac{x^2-y^2}{(x^2+y^2)^2}$ で、たすと0（原点以外）。点 $(1,\ 2)$ では $3\times\dfrac{4-1}{25}=0.36$ と $3\times\dfrac{1-4}{25}=-0.36$。}

\begin{furikaeri}
\begin{checklist}
\item 矢印の図にマスを置いて、わき出しと吸いこみを見分けられる。\bkfurref{bk:8-9:hako}
\item マスの辺を通る量から、表計算で発散を求められる。\bkfurref{bk:8-9:masu}
\item 発散を偏微分の和で計算し、流束や連続の式と結びつけられる。\bkfurref{bk:8-9:shiki}
\end{checklist}
\end{furikaeri}

\bkfukidashi{mina}{矢印を見ているだけではわからなかった池の口が、たし算とひき算で見つかったね。次は、矢印が「回る」ようすを調べよう。}
